% ---------------------------------------------------------------------------
% Methodology: MADAR + Unlearning
%
% SECTION BODY ONLY -- no \documentclass, no preamble. Drop into a paper with
%   % ---------------------------------------------------------------------------
% Methodology: MADAR + Unlearning
%
% SECTION BODY ONLY -- no \documentclass, no preamble. Drop into a paper with
%   % ---------------------------------------------------------------------------
% Methodology: MADAR + Unlearning
%
% SECTION BODY ONLY -- no \documentclass, no preamble. Drop into a paper with
%   % ---------------------------------------------------------------------------
% Methodology: MADAR + Unlearning
%
% SECTION BODY ONLY -- no \documentclass, no preamble. Drop into a paper with
%   \input{docs/methodology_madar_unlearning}
% or compile the preview wrapper docs/methodology_standalone.tex.
%
% Required packages (see the wrapper for a working preamble):
%   amsmath, amssymb, booktabs, algorithm, algpseudocode, xcolor, hyperref
%
% Every equation, constant and default below is transcribed from the
% implementation in code/; the trailing comment on a line names the source.
% ---------------------------------------------------------------------------

\section{Methodology}
\label{sec:method}

\subsection{Problem setting and notation}
\label{sec:method:notation}

We study class-incremental learning over a stream of $T{+}1$ tasks. A task
setup is written $C_0 + s \times T$: task $0$ introduces $C_0$ base families and
each subsequent task $t \in \{1,\dots,T\}$ introduces $s$ previously unseen
families, for $K = C_0 + sT$ classes in total.\footnote{\texttt{core/tasks.py}.
The EMBER protocol is $50+5\times10$ ($K=100$); LAMDA is $30+5\times10$
($K=80$).} Class indices follow \emph{presentation order}, so the classes
introduced by task $t$ are
\begin{equation}
  \mathcal{C}_t =
  \begin{cases}
    \{0,\dots,C_0-1\}, & t = 0,\\[2pt]
    \{C_0 + (t-1)s, \dots, C_0 + ts - 1\}, & t \ge 1,
  \end{cases}
\end{equation}
and the \emph{active prefix} after task $t$ has width $a_t = C_0 + ts$, with
$a_{-1} \coloneqq C_0$ by convention. Ordering the classes this way is what makes
prefix masking well defined: it is not a cosmetic relabelling.

Let $f_\theta : \mathcal{X} \to \mathbb{R}^{K}$ be the classifier, emitting a
logit vector of full width $K$ at every stage, and let
$h_\theta : \mathcal{X} \to \mathbb{R}^{d_{\text{lat}}}$ be its penultimate
(latent) representation. Classes that have not yet arrived are suppressed by an
additive mask rather than by reshaping the head,
\begin{equation}
  \label{eq:mask}
  m^{(a)} \in \mathbb{R}^{K}, \qquad
  m^{(a)}_j = \begin{cases} 0, & j < a, \\ -10^{9}, & j \ge a, \end{cases}
\end{equation}
so a single parameterisation carries the whole stream and no logit belonging to
an unseen class can contribute to a loss or to an $\arg\max$.%
\footnote{\texttt{core/training.py:logit\_mask}.}
Table~\ref{tab:notation} collects the remaining symbols.

\begin{table}[t]
  \centering
  \small
  \caption{Notation, and where each quantity is realised in the implementation.}
  \label{tab:notation}
  \begin{tabular}{@{}llp{0.30\linewidth}@{}}
    \toprule
    Symbol & Meaning & Implementation \\
    \midrule
    $\mathcal{D}_t$ & training samples of task $t$ & \texttt{Experiment.task\_loader} \\
    $a_t$ & active logit width after task $t$ & \texttt{TaskSchedule.active\_count} \\
    $m^{(a)}$ & additive prefix mask, Eq.~\eqref{eq:mask} & \texttt{training.logit\_mask} \\
    $\mathcal{M}_t$ & replay buffer after task $t$ & \texttt{core.buffer.ReplayBuffer} \\
    $M$ & total buffer capacity & \texttt{mem\_size} \\
    $\theta^{-}$ & frozen teacher & \texttt{training.clone\_teacher} \\
    $\theta^{\text{pre}}_t$ & model before the forget phase & \texttt{pre} in \texttt{train\_task} \\
    $\mathcal{F}_t,\ \mathcal{R}_t$ & forget / retain split of $\mathcal{D}_t$ & \texttt{select\_forget\_set} \\
    $\rho$ & forget budget as a fraction of $|\mathcal{D}_t|$ & \texttt{forget\_ratio} \\
    $\alpha$ & weight on the forget term & \texttt{alpha} \\
    $\gamma_t$ & weight on the current-task term & \texttt{training.rnt\_for} \\
    $\Omega,\ \theta^{\text{old}}$ & SI importances and anchor & \texttt{SynapticIntelligence} \\
    $c$ & SI strength & \texttt{si\_c} \\
    $\tau$ & distillation temperature & \texttt{kd\_temp} \\
    \bottomrule
  \end{tabular}
\end{table}

\paragraph{Preprocessing.}
Families enter the study only if they carry at least $200$ training samples, and
the $K$ most frequent eligible families are retained. Feature standardisation is
fit \emph{on task~0 alone} and applied to the whole stream, clipped to
$\pm 10$ standard deviations; fitting on the full training set would leak the
statistics of families the learner has not yet met into a protocol whose entire
claim is that it has not met them. LAMDA's features are binary and are left
unscaled.

\subsection{Shared continual backbone}
\label{sec:method:madar}

All conditions share one training loop, one evaluation protocol and one
determinism block, and differ only in the rule applied to a task
(\texttt{experiments/base.py}). We describe that shared backbone first, because
MADAR${+}$Unlearning is defined as a strict extension of it: everything in this
subsection is common to both, so any difference measured between them is
attributable to Section~\ref{sec:method:unlearn} alone.

\paragraph{Task 0.}
The base families are fit by ordinary supervised training for $E_0$ epochs with
Adam,
\begin{equation}
  \label{eq:task0}
  \min_\theta\ \mathbb{E}_{(x,y)\sim\mathcal{D}_0}
    \big[\ \ell_{\mathrm{CE}}\!\big(f_\theta(x) + m^{(a_0)},\, y\big)\ \big].
\end{equation}
Afterwards the SI anchor is set, $\theta^{\text{old}} \leftarrow \theta$, the
teacher is cloned, $\theta^{-} \leftarrow \theta$, and the buffer is filled from
$\mathcal{D}_0$ by the rule of Section~\ref{sec:method:buffer}.

\paragraph{Continual phase.}
For each $t \ge 1$ the model is trained for a \emph{fixed iteration budget}
$N_{\mathrm{cl}}$ rather than a fixed number of epochs, so that compute per task
does not track task size. Each step draws one batch $(x,y) \sim \mathcal{D}_t$
and one replay batch $(\tilde{x},\tilde{y}) \sim \mathcal{M}_{t-1}$ and minimises
\begin{align}
  \label{eq:cl}
  \mathcal{L}^{\mathrm{CL}}_t
   &= \gamma_t\,
      \ell_{\mathrm{CE}}\!\big(f_\theta(x \cup \tilde{x}) + m^{(a_t)},\;
                                y \cup \tilde{y}\big) \nonumber\\
   &\quad + (1-\gamma_t)\,
      \mathcal{L}_{\mathrm{KD}}\!\big(\tilde{x};\, \theta,\, \theta^{-},\, a_{t-1}\big)
      \;+\; c\,\Omega(\theta),
\end{align}
where the distillation term is the temperature-scaled divergence \emph{from} the
teacher \emph{to} the student, restricted to the previous active prefix,
\begin{equation}
  \label{eq:kd}
  \mathcal{L}_{\mathrm{KD}}(x;\theta,\theta^{-},a)
   = \tau^{2}\,
     \mathrm{KL}\!\Big(
       \sigma\big(f_{\theta^{-}}(x)_{:a}/\tau\big)
       \,\Big\|\,
       \sigma\big(f_{\theta}(x)_{:a}/\tau\big)
     \Big),
\end{equation}
with $\sigma$ the softmax. Distillation is taken over \emph{replayed} samples
only, which is why the implementation refuses the configuration
$\texttt{use\_kd}\wedge\neg\texttt{use\_replay}$ outright rather than silently
distilling on the current task.

The convex weight is $\gamma_t = 1/(t+1)$, which weights every task seen so far
equally. Three alternatives (\texttt{floor}, \texttt{sqrt}, \texttt{fixed}) are
implemented and default to off; they exist because $1/(t+1)$ is principled only
when tasks are interchangeable, and under a task~0 far larger than the
increments the newest families receive a small share of the gradient. We note
also that $\gamma_t$ is meaningful only \emph{as} a convex weight against the KD
term: in an ablation that removes distillation, retaining $\gamma_t$ degenerates
into a bare $1/(t+1)$ scalar on the whole objective --- a decaying learning rate
in the costume of a replay weight. It is therefore disabled by default whenever
KD is disabled.

\paragraph{Synaptic Intelligence.}
Parameter importance is the path integral of Zenke et al. During training, each
optimiser step contributes
\begin{equation}
  \label{eq:si-accum}
  W_k \;\mathrel{+}=\; -\,g_k \cdot \Delta\theta_k,
\end{equation}
with $g_k$ the pre-step gradient and $\Delta\theta_k$ the realised parameter
change. At a task boundary the running sum is folded into the importances,
normalised by the squared distance the parameter actually travelled,
\begin{equation}
  \label{eq:si-omega}
  \Omega_k \;\mathrel{+}=\;
    \frac{W_k}{\big(\theta_k - \theta^{\text{old}}_k\big)^{2} + \xi},
  \qquad W_k \leftarrow 0,
\end{equation}
and the quadratic penalty charged during subsequent training is
\begin{equation}
  \label{eq:si-penalty}
  \Omega(\theta) \;=\; \sum_k \Omega_k \big(\theta_k - \theta^{\text{old}}_k\big)^{2}.
\end{equation}
Whether the anchor $\theta^{\text{old}}$ advances at the boundary is a separate
decision from whether $\Omega$ is updated, and Section~\ref{sec:method:order}
turns on exactly that separation.

\paragraph{Frozen normalisation statistics.}
BatchNorm layers are held in evaluation mode for the entire continual phase.
Their running statistics were estimated on task~0, and allowing a stream
dominated by the newest families to overwrite them is a channel of forgetting
that is unrelated to the weights and would otherwise be silently attributed to
the training rule.

\subsection{Diversity-aware replay buffer}
\label{sec:method:buffer}

The buffer holds at most $M$ samples and is budgeted uniformly across the
families seen so far: with $F$ distinct families held or arriving, the
per-family budget is $b = \max(1, \lfloor M / F \rfloor)$, existing families are
truncated to $b$, and each arriving family is filled to $b$.

Within a family, selection is MADAR's diversity criterion. An Isolation Forest
is fitted over the family's samples and they are ordered ascending by
$\mathrm{decision\_function}$, i.e. most anomalous first:
\begin{equation}
  \label{eq:buffer-order}
  \pi = \operatorname{argsort}\big(\mathrm{IF}(\,\cdot\,)\big),
  \qquad
  \mathcal{M}^{(f)} = \operatorname{interleave}\big(
     \pi_{1:\lfloor b/2 \rfloor},\;
     \pi_{n - \lceil b/2 \rceil + 1 : n}\big),
\end{equation}
keeping the most anomalous half and the most prototypical half. The two ends are
\emph{interleaved} rather than concatenated so that a later truncation to a
smaller budget --- which takes a prefix of the stored list --- remains balanced
between the two ends instead of retaining whichever end happens to sit first.

Two spaces are selectable for the forest. Fitting on the scaled \emph{input}
features (\texttt{--buffer-space raw}) is MADAR as published: selection is then a
fixed property of the data. Fitting on the classifier's latent representation
$h_\theta$ (\texttt{latent}, the default here; MADAR-$\theta$) makes selection
move as the model learns. The distinction is not cosmetic --- it decides whether
the memory is chosen in a fixed space or in one the learner is simultaneously
reshaping --- and it is recorded per run.\footnote{The Isolation Forest's
\texttt{contamination} parameter appears in the configuration but only shifts the
estimator's decision \emph{threshold}. Selection here is purely rank-based on
$\mathrm{decision\_function}$ scores, so its value does not change which samples
are chosen; it is retained because the estimator requires it.}

\subsection{MADAR + Unlearning}
\label{sec:method:unlearn}

The proposed condition is the backbone of Section~\ref{sec:method:madar} with a
\emph{forget phase} appended to every task $t \ge 1$. A fraction of the task the
model has just learned is nominated for deletion, the model is trained to become
uninformative about it, and the replay buffer is rebuilt so that the deleted
samples cannot return. The phase runs \emph{after} the continual phase, so the
setting is learn-then-forget: the samples in question have already been fit.
Algorithm~\ref{alg:madaru} states the whole procedure; the subsections that
follow justify the parts that are not mechanical.

\begin{algorithm}[t]
\caption{MADAR${+}$Unlearning, one task $t \ge 1$}
\label{alg:madaru}
\begin{algorithmic}[1]
  \Require model $\theta$, teacher $\theta^{-}$, buffer $\mathcal{M}_{t-1}$,
           SI state $(W,\Omega,\theta^{\text{old}})$
  \State $\theta \gets \textsc{ContinualPhase}(\theta,\theta^{-},\mathcal{D}_t,\mathcal{M}_{t-1})$
         \Comment{Eq.~\eqref{eq:cl}, $N_{\mathrm{cl}}$ steps}
  \State $\Omega \gets \textsc{FoldImportances}(W,\theta,\theta^{\text{old}})$
         \Comment{Eq.~\eqref{eq:si-omega}; anchor \emph{not} advanced}
  \State $(\mathcal{F}_t, \mathcal{R}_t) \gets \textsc{Select}(\mathcal{D}_t, \rho)$
         \Comment{Section~\ref{sec:method:select}}
  \If{$\mathcal{F}_t = \emptyset$}
      \State advance anchor, refresh teacher, rebuild buffer from $\mathcal{D}_t$;
             \textbf{return}
  \EndIf
  \State $\theta^{\text{pre}} \gets \theta$ \Comment{frozen; teacher and reference}
  \For{$e = 1$ \textbf{to} $E_u$}
    \For{minibatch $x_f \subset \mathcal{F}_t$}
      \State draw $x_r \sim \mathcal{R}_t$, \; $\tilde{x} \sim \mathcal{M}_{t-1}$
      \State $\theta \gets \theta - \eta_u \nabla_\theta \mathcal{L}^{U}$
             \Comment{Eq.~\eqref{eq:unlearn}}
    \EndFor
  \EndFor
  \State $\textsc{Measure}(\theta^{\text{pre}}, \theta, \mathcal{F}_t, \mathcal{R}_t)$
         \Comment{Section~\ref{sec:method:efficacy}}
  \State $\theta^{\text{old}} \gets \theta$;\quad $\theta^{-} \gets \theta$
  \State $\mathcal{M}_t \gets \textsc{RebuildBuffer}(\mathcal{M}_{t-1}, \mathcal{R}_t)$
         \Comment{retain split only}
\end{algorithmic}
\end{algorithm}

\subsubsection{Ordering of the SI boundary}
\label{sec:method:order}

The task boundary is closed \emph{before} the forget phase, but the anchor is
deliberately held back: Eq.~\eqref{eq:si-omega} is applied with
$\theta^{\text{old}}$ left pointing at the parameters as they stood at the
\emph{start} of task $t$, and it advances only once unlearning is complete
(lines 2 and 13 of Algorithm~\ref{alg:madaru}).

This ordering is load-bearing in both directions. Folding $W$ into $\Omega$
first means the importances already account for what task $t$ made important, so
the forget phase is constrained by a current estimate rather than a stale one.
Withholding the anchor means the penalty in Eq.~\eqref{eq:si-penalty} is
evaluated against the start of the task, so unlearning is charged for every
important parameter it drags away from that point. Had the anchor advanced
first, $\theta - \theta^{\text{old}} = 0$ at the moment the forget phase begins
and the regulariser would be identically inert for precisely the phase it is
meant to regulate.

We note for completeness that the path integral is \emph{not} accumulated during
the forget phase --- Eq.~\eqref{eq:si-accum} is applied only inside the
continual phase, since $W$ has already been consumed at line~2. The unlearning
trajectory is therefore regularised by SI but does not itself contribute to
future importance estimates.

\subsubsection{Forget-set selection}
\label{sec:method:select}

Selection partitions the current task's samples into
$\mathcal{D}_t = \mathcal{F}_t \sqcup \mathcal{R}_t$ under a budget
$k = \lfloor \rho\,|\mathcal{D}_t| \rfloor$ that is held \emph{fixed across
selectors}. Holding the budget fixed is what makes a comparison between rules
attributable to \emph{which} samples are named rather than to how many. The
seed for the decision is derived, $\mathrm{seed}_t = \mathrm{child}(\mathrm{seed}, t, 7)$,
rather than drawn from the global stream, so that introducing a selector cannot
shift every later random draw in the run and thereby de-pair two conditions
intended to share a seed. If $k = 0$ the phase is skipped and the run record
says so explicitly.

Three rules are compared.

\begin{description}
  \item[\texttt{random}] A uniform subset of size $k$. This is the control for
    \emph{which} samples: it isolates the effect of the selection rule from the
    effect of deleting $k$ samples at all.

  \item[\texttt{donut}] (default) An Isolation Forest is fitted jointly over the
    task in latent space, samples are ordered ascending by
    $\mathrm{decision\_function}$, and the forget set is the \emph{middle band}
    of that ordering,
    \begin{equation}
      \label{eq:donut}
      \mathcal{F}_t = \pi_{\,j:\,j+k},
      \qquad j = \Big\lfloor \tfrac{n}{2} \Big\rfloor - \Big\lfloor \tfrac{k}{2} \Big\rfloor .
    \end{equation}
    Both the most anomalous and the most prototypical samples are kept. This
    applies MADAR's own prior --- the criterion it uses to \emph{choose} what to
    remember, Eq.~\eqref{eq:buffer-order} --- to the complementary question of
    what to delete, and is the hand-crafted rule the learned policy is later
    measured against.

  \item[\texttt{leftover}] A per-family Isolation Forest in \emph{raw} feature
    space; each family keeps the anomalous/prototypical halves that a
    MADAR-style selection at the current per-family budget would have kept, and
    everything else is forgotten. This rule ignores $\rho$ by construction: the
    forget set is whatever the buffer budget leaves over, which is typically far
    larger than a $10\%$ slab. Comparisons involving it are therefore not
    budget-matched, and we report it as such.
\end{description}

\subsubsection{The forget objective}
\label{sec:method:objective}

Let $\theta^{\text{pre}}$ denote the model at the end of the continual phase,
frozen and used both as the distillation teacher for this phase and as the
reference for efficacy measurement. The retain anchor is therefore ``what the
model believed a moment ago'', not the previous task's model. Writing
$u = \mathbf{1}/s$ for the uniform distribution over the $s = a_t - a_{t-1}$
classes this task introduced, the phase minimises
\begin{align}
  \label{eq:unlearn}
  \mathcal{L}^{U}
   &= \alpha\,
      \underbrace{\mathrm{KL}\!\Big(u \,\Big\|\,
        \sigma\big(f_\theta(x_f)_{a_{t-1}:a_t}\big)\Big)}_{\text{forget}}
      \nonumber\\
   &\quad + (1-\alpha)\,
      \underbrace{\Big[
        \tfrac{1}{2}\,\ell_{\mathrm{CE}}\big(f_\theta(x_\cup) + m^{(a_t)},\, y_\cup\big)
        + \tfrac{1}{2}\,\mathcal{L}_{\mathrm{KD}}\big(x_\cup; \theta, \theta^{\text{pre}}, a_t\big)
      \Big]}_{\text{retain}}
      \nonumber\\
   &\quad + c\,\Omega(\theta),
\end{align}
where $x_f \subset \mathcal{F}_t$ and
$(x_\cup, y_\cup) = (x_r,y_r) \cup (\tilde{x},\tilde{y})$ concatenates a batch of
this task's \emph{retained} samples with a replay batch from the buffer. The
forget set defines the step count; the retain and replay streams cycle to match
it. Optimisation uses Adam at $\eta_u$ for $E_u$ passes over $\mathcal{F}_t$,
with gradient-norm clipping at $0.5$ and BatchNorm still frozen.

Three properties of Eq.~\eqref{eq:unlearn} should be read carefully.

\emph{The forget term acts on the new-task slice only.} The softmax is taken over
$[a_{t-1}, a_t)$ --- the $s$ classes this task introduced --- not over the full
active prefix. The objective therefore drives the model toward indifference
\emph{among the new families}, which is a narrower claim than indifference among
all classes seen so far.

\emph{Retention is anchored on two populations.} Because $x_\cup$ mixes retained
current-task samples with replayed old-family samples, the $(1-\alpha)$ term
defends both plasticity on the task being partially forgotten and stability on
everything before it, under one weight.

\emph{The regulariser is not slack.} By Section~\ref{sec:method:order}, $\Omega$
is evaluated against the start of the task, so the forget term must overcome an
active quadratic penalty on exactly those parameters the task made important.

\subsubsection{Buffer reconstruction from the retain split}
\label{sec:method:rebuild}

After the forget phase the buffer is rebuilt from $\mathcal{R}_t$ rather than
$\mathcal{D}_t$: latents are recomputed with the \emph{post}-unlearning model and
Eq.~\eqref{eq:buffer-order} is re-applied per family. Forgotten samples are
consequently never eligible for replay at any later task.

This is worth separating from the objective, because it is a second and
independent deletion mechanism, and it operates whenever a forget set is selected
--- including when $\alpha = 0$. It is the reason the ablation of
Section~\ref{sec:method:controls} is necessary rather than decorative.

\subsubsection{Efficacy measurement}
\label{sec:method:efficacy}

Unlearning is measured, per task, by comparing $\theta^{\text{pre}}$ against the
post-phase $\theta$ on the same inputs, over two disjoint scopes: the forget set
$\mathcal{F}_t$ and the retained current-task samples $\mathcal{R}_t$. For each
scope we record accuracy before and after, the mean predictive entropy and its
normalised form, and the displacement of the latent representation:
\begin{align}
  \label{eq:efficacy}
  H(x) &= -\!\!\sum_{j < a_t} p_j \log p_j,
  \qquad p = \sigma\big(f_{\theta}(x)_{:a_t}\big), \\
  \widehat{H} &= \frac{\mathbb{E}[H(x)]}{\log a_t},
  \qquad
  \Delta_{\text{lat}}(x) = \big\| h_{\theta}(x) - h_{\theta^{\text{pre}}}(x) \big\|_2 .
\end{align}
A useful result is a large drop in accuracy and a large $\Delta_{\text{lat}}$ on
$\mathcal{F}_t$ together with small movement on $\mathcal{R}_t$: forgetting that
also flattens the retained samples is collateral damage, not deletion. Note that
$\widehat{H}$ is normalised over the \emph{full} active prefix while the forget
term of Eq.~\eqref{eq:unlearn} operates on the $s$-class slice, so a forget set
driven to perfect uniformity within that slice will not register
$\widehat{H} \approx 1$; the two widths are not interchangeable and we report
$\widehat{H}$ as a monotone diagnostic rather than as a saturating one.

\subsection{Conditions, controls and ablations}
\label{sec:method:controls}

Four conditions share the network, the family selection, the scaler, the task
schedule, the masked evaluation and the determinism block. Only the rule applied
to a task differs.

\begin{table}[t]
  \centering
  \small
  \caption{The four conditions. All share every non-training component.}
  \label{tab:conditions}
  \begin{tabular}{@{}lp{0.62\linewidth}@{}}
    \toprule
    Condition & Training rule \\
    \midrule
    Naive & Fine-tune on the current task only. Lower bound. \\
    Joint & Discard the model and retrain from scratch on all data seen so far.
            Oracle upper bound; not a realisable continual setup. \\
    MADAR & Eq.~\eqref{eq:cl}: diversity-aware replay, distillation against the
            previous task's model, SI penalty. \\
    MADAR${+}$U & MADAR, plus the per-task forget phase of
            Section~\ref{sec:method:unlearn}. \\
    \bottomrule
  \end{tabular}
\end{table}

\paragraph{$\alpha = 0$ is the matched control.}
The headline comparison is not MADAR${+}$U against MADAR. At $\alpha = 0$ the
forget set is still selected by the same rule, is still excluded from the buffer
rebuild, and the same number of gradient steps is still taken --- only the push
toward uniform is absent. The contrast $\alpha > 0$ versus $\alpha = 0$ therefore
isolates the value of the \emph{forget objective}, while $\alpha = 0$ versus
MADAR isolates the value of \emph{curation}, i.e. of removing those samples from
memory at all. Reporting the method without the $\alpha = 0$ arm would attribute
to unlearning what may simply be memory curation, and both arms are run for every
corpus and seed.

\paragraph{Component ablations.}
Replay, distillation and the SI penalty are independently switchable in
Eq.~\eqref{eq:cl} so that an ablation removes exactly one thing. SI is ablated by
setting $c = 0$, which leaves the optimiser, the iteration count and the data
order untouched; removing distillation additionally disables $\gamma_t$ for the
reason given in Section~\ref{sec:method:madar}.

\subsection{Evaluation protocol}
\label{sec:method:eval}

After every task, predictions are taken as
$\hat{y} = \arg\max_{j < a_t} f_\theta(x)_j$ over the test split, and all
reported quantities are derived from a single confusion matrix so that the
numbers in a record are mutually consistent by construction rather than by three
code paths agreeing. Three scopes are evaluated:
\begin{itemize}
  \item \textbf{seen} --- all families encountered so far. The headline number.
  \item \textbf{recent} --- the families just introduced. Acquisition, i.e.
        plasticity.
  \item \textbf{task0} --- the base families. Retention, i.e. resistance to
        forgetting.
\end{itemize}
The \textbf{seen} scope averages the other two, so a method that trades
acquisition against retention is invisible in it alone; we report all three.

Macro averages are taken over classes \emph{with support in the evaluation
slice}. A class absent from a slice has no recall to measure, and scoring it as
zero would make the metric a function of which classes happen to be absent.
Because a class may be predicted without being present, the false positives such
predictions generate belong to no scored class; they are surfaced explicitly as
\texttt{n\_pred\_outside\_support} rather than dropped. For single-label
multiclass, micro-precision, micro-recall and micro-F1 all equal accuracy, so
only accuracy is reported under that name.

Two protocol constraints are enforced rather than documented. First, the
\texttt{stage1} protocol (task~0 retrained per seed) and the \texttt{meta}
protocol (a shared task-0 checkpoint) are never differenced against each other,
and are tabulated as separate blocks. Second, a crashed run is recorded as failed
and is never scored, defaulted or interpolated --- a missing number that silently
becomes \texttt{NaN} and then sorts to the top of a ranking is a real failure
mode, and the sweep boundary is where it is refused.

\subsection{Hyperparameters}
\label{sec:method:hparams}

Table~\ref{tab:hparams} gives the values in effect. Per-corpus entries are
protocol for that corpus; an explicit override always wins and, whatever wins,
the resolved value is written into the run record --- there is no value in effect
that the log does not name.

\begin{table}[t]
  \centering
  \small
  \caption{Protocol parameters. Shared values are common to every corpus.}
  \label{tab:hparams}
  \begin{tabular}{@{}llcccc@{}}
    \toprule
    & & EMBER 2018 & EMBER 2024 & LAMDA & Tiny ImageNet \\
    \midrule
    \multicolumn{6}{@{}l}{\emph{Shared}}\\
    Task-0 epochs $E_0$        & \texttt{task0\_epochs}   & \multicolumn{3}{c}{30, Adam $10^{-3}$} & 50 \\
    Batch size                 & \texttt{batch\_size}     & \multicolumn{3}{c}{256} & 32 \\
    Continual optimiser        & ---                      & \multicolumn{4}{c}{SGD, momentum $0.9$, weight decay $10^{-6}$, clip $1.0$} \\
    Distillation temperature   & \texttt{kd\_temp}        & \multicolumn{4}{c}{$\tau = 2.0$} \\
    Current-task weight        & \texttt{rnt\_mode}       & \multicolumn{4}{c}{$\gamma_t = 1/(t{+}1)$} \\
    Forget weight              & \texttt{alpha}           & \multicolumn{4}{c}{$\alpha \in \{0.2,\ 0\}$ (both arms)} \\
    Forget budget              & \texttt{forget\_ratio}   & \multicolumn{4}{c}{$\rho = 0.10$} \\
    Forget passes              & \texttt{unlearn\_epochs} & \multicolumn{3}{c}{$E_u = 3$} & 3 \\
    Forget optimiser           & \texttt{unlearn\_lr}     & \multicolumn{4}{c}{Adam $\eta_u = 10^{-4}$, clip $0.5$} \\
    Family eligibility         & \texttt{min\_family\_samples} & \multicolumn{3}{c}{$\ge 200$} & --- \\
    \midrule
    \multicolumn{6}{@{}l}{\emph{Per corpus}}\\
    Schedule                   & \texttt{--task-setup}    & $50{+}5{\times}10$ & $50{+}5{\times}10$ & $30{+}5{\times}10$ & $100{+}10{\times}10$ \\
    Continual budget $N_{\mathrm{cl}}$ & \texttt{cl\_iters} & 2000 & 8000 & 4000 & 3000 \\
    Continual lr               & \texttt{cl\_lr}          & $10^{-4}$ & $10^{-4}$ & $10^{-4}$ & $0.03$ \\
    Buffer capacity $M$        & \texttt{mem\_size}       & 5000 & 10000 & 2500 & 2000 \\
    SI strength $c$            & \texttt{si\_c}           & 1.0 & 100.0 & 100.0 & 1.0 \\
    Per-family cap             & \texttt{family\_cap}     & --- & 10000 & --- & --- \\
    Feature clip               & \texttt{feature\_clip}   & $\pm10$ & $\pm10$ & --- (binary) & --- \\
    \bottomrule
  \end{tabular}
\end{table}

\subsection{Threats to validity}
\label{sec:method:threats}

\paragraph{Compute is not matched between MADAR and MADAR${+}$U.}
The forget phase adds $E_u \lceil |\mathcal{F}_t| / B \rceil$ optimiser steps per
task on top of the $N_{\mathrm{cl}}$ continual steps, and the run record reports
the total. The difference is small relative to $N_{\mathrm{cl}}$ but it is not
zero, so a MADAR${+}$U improvement over MADAR is confounded with additional
gradient steps. The $\alpha = 0$ arm of Section~\ref{sec:method:controls} is
matched in step count and is the comparison from which the forget objective's
contribution should be read.

\paragraph{The forget objective and the efficacy metric span different widths.}
As noted in Section~\ref{sec:method:efficacy}, $\mathcal{L}^U$ flattens the
$s$-class new-task slice while $\widehat{H}$ normalises over the full active
prefix of width $a_t$. The two are consistent in direction but not on a common
scale.

\paragraph{Unlearning is behavioural, not certified.}
Eq.~\eqref{eq:unlearn} drives the model's \emph{predictions} on $\mathcal{F}_t$
toward uniformity and removes those samples from memory; it offers no guarantee
of indistinguishability from a model retrained without them, and we make no such
claim. The reported quantities of Section~\ref{sec:method:efficacy} are evidence
of behavioural forgetting under a specific probe, and the \texttt{Joint} oracle
of Table~\ref{tab:conditions} is a retraining reference for accuracy, not a
privacy reference.

\paragraph{The \texttt{leftover} selector is not budget-matched.}
Its forget set is determined by the buffer budget rather than by $\rho$, so
differences against \texttt{random} and \texttt{donut} confound the selection
rule with the number of samples deleted.

\paragraph{Reproduction.}
Every reported run is produced by a single command and writes one self-describing
JSON record naming the corpus, the schedule, the model and parameter count, every
hyperparameter in effect, the seed, the protocol, the library versions, the GPU
and the git commit:
\begin{verbatim}
./iclr run --dataset ember2018 --task-setup ember \
           --experiment madar_unlearn --selector donut \
           --model full --seed 1 --group <group>
\end{verbatim}

% or compile the preview wrapper docs/methodology_standalone.tex.
%
% Required packages (see the wrapper for a working preamble):
%   amsmath, amssymb, booktabs, algorithm, algpseudocode, xcolor, hyperref
%
% Every equation, constant and default below is transcribed from the
% implementation in code/; the trailing comment on a line names the source.
% ---------------------------------------------------------------------------

\section{Methodology}
\label{sec:method}

\subsection{Problem setting and notation}
\label{sec:method:notation}

We study class-incremental learning over a stream of $T{+}1$ tasks. A task
setup is written $C_0 + s \times T$: task $0$ introduces $C_0$ base families and
each subsequent task $t \in \{1,\dots,T\}$ introduces $s$ previously unseen
families, for $K = C_0 + sT$ classes in total.\footnote{\texttt{core/tasks.py}.
The EMBER protocol is $50+5\times10$ ($K=100$); LAMDA is $30+5\times10$
($K=80$).} Class indices follow \emph{presentation order}, so the classes
introduced by task $t$ are
\begin{equation}
  \mathcal{C}_t =
  \begin{cases}
    \{0,\dots,C_0-1\}, & t = 0,\\[2pt]
    \{C_0 + (t-1)s, \dots, C_0 + ts - 1\}, & t \ge 1,
  \end{cases}
\end{equation}
and the \emph{active prefix} after task $t$ has width $a_t = C_0 + ts$, with
$a_{-1} \coloneqq C_0$ by convention. Ordering the classes this way is what makes
prefix masking well defined: it is not a cosmetic relabelling.

Let $f_\theta : \mathcal{X} \to \mathbb{R}^{K}$ be the classifier, emitting a
logit vector of full width $K$ at every stage, and let
$h_\theta : \mathcal{X} \to \mathbb{R}^{d_{\text{lat}}}$ be its penultimate
(latent) representation. Classes that have not yet arrived are suppressed by an
additive mask rather than by reshaping the head,
\begin{equation}
  \label{eq:mask}
  m^{(a)} \in \mathbb{R}^{K}, \qquad
  m^{(a)}_j = \begin{cases} 0, & j < a, \\ -10^{9}, & j \ge a, \end{cases}
\end{equation}
so a single parameterisation carries the whole stream and no logit belonging to
an unseen class can contribute to a loss or to an $\arg\max$.%
\footnote{\texttt{core/training.py:logit\_mask}.}
Table~\ref{tab:notation} collects the remaining symbols.

\begin{table}[t]
  \centering
  \small
  \caption{Notation, and where each quantity is realised in the implementation.}
  \label{tab:notation}
  \begin{tabular}{@{}llp{0.30\linewidth}@{}}
    \toprule
    Symbol & Meaning & Implementation \\
    \midrule
    $\mathcal{D}_t$ & training samples of task $t$ & \texttt{Experiment.task\_loader} \\
    $a_t$ & active logit width after task $t$ & \texttt{TaskSchedule.active\_count} \\
    $m^{(a)}$ & additive prefix mask, Eq.~\eqref{eq:mask} & \texttt{training.logit\_mask} \\
    $\mathcal{M}_t$ & replay buffer after task $t$ & \texttt{core.buffer.ReplayBuffer} \\
    $M$ & total buffer capacity & \texttt{mem\_size} \\
    $\theta^{-}$ & frozen teacher & \texttt{training.clone\_teacher} \\
    $\theta^{\text{pre}}_t$ & model before the forget phase & \texttt{pre} in \texttt{train\_task} \\
    $\mathcal{F}_t,\ \mathcal{R}_t$ & forget / retain split of $\mathcal{D}_t$ & \texttt{select\_forget\_set} \\
    $\rho$ & forget budget as a fraction of $|\mathcal{D}_t|$ & \texttt{forget\_ratio} \\
    $\alpha$ & weight on the forget term & \texttt{alpha} \\
    $\gamma_t$ & weight on the current-task term & \texttt{training.rnt\_for} \\
    $\Omega,\ \theta^{\text{old}}$ & SI importances and anchor & \texttt{SynapticIntelligence} \\
    $c$ & SI strength & \texttt{si\_c} \\
    $\tau$ & distillation temperature & \texttt{kd\_temp} \\
    \bottomrule
  \end{tabular}
\end{table}

\paragraph{Preprocessing.}
Families enter the study only if they carry at least $200$ training samples, and
the $K$ most frequent eligible families are retained. Feature standardisation is
fit \emph{on task~0 alone} and applied to the whole stream, clipped to
$\pm 10$ standard deviations; fitting on the full training set would leak the
statistics of families the learner has not yet met into a protocol whose entire
claim is that it has not met them. LAMDA's features are binary and are left
unscaled.

\subsection{Shared continual backbone}
\label{sec:method:madar}

All conditions share one training loop, one evaluation protocol and one
determinism block, and differ only in the rule applied to a task
(\texttt{experiments/base.py}). We describe that shared backbone first, because
MADAR${+}$Unlearning is defined as a strict extension of it: everything in this
subsection is common to both, so any difference measured between them is
attributable to Section~\ref{sec:method:unlearn} alone.

\paragraph{Task 0.}
The base families are fit by ordinary supervised training for $E_0$ epochs with
Adam,
\begin{equation}
  \label{eq:task0}
  \min_\theta\ \mathbb{E}_{(x,y)\sim\mathcal{D}_0}
    \big[\ \ell_{\mathrm{CE}}\!\big(f_\theta(x) + m^{(a_0)},\, y\big)\ \big].
\end{equation}
Afterwards the SI anchor is set, $\theta^{\text{old}} \leftarrow \theta$, the
teacher is cloned, $\theta^{-} \leftarrow \theta$, and the buffer is filled from
$\mathcal{D}_0$ by the rule of Section~\ref{sec:method:buffer}.

\paragraph{Continual phase.}
For each $t \ge 1$ the model is trained for a \emph{fixed iteration budget}
$N_{\mathrm{cl}}$ rather than a fixed number of epochs, so that compute per task
does not track task size. Each step draws one batch $(x,y) \sim \mathcal{D}_t$
and one replay batch $(\tilde{x},\tilde{y}) \sim \mathcal{M}_{t-1}$ and minimises
\begin{align}
  \label{eq:cl}
  \mathcal{L}^{\mathrm{CL}}_t
   &= \gamma_t\,
      \ell_{\mathrm{CE}}\!\big(f_\theta(x \cup \tilde{x}) + m^{(a_t)},\;
                                y \cup \tilde{y}\big) \nonumber\\
   &\quad + (1-\gamma_t)\,
      \mathcal{L}_{\mathrm{KD}}\!\big(\tilde{x};\, \theta,\, \theta^{-},\, a_{t-1}\big)
      \;+\; c\,\Omega(\theta),
\end{align}
where the distillation term is the temperature-scaled divergence \emph{from} the
teacher \emph{to} the student, restricted to the previous active prefix,
\begin{equation}
  \label{eq:kd}
  \mathcal{L}_{\mathrm{KD}}(x;\theta,\theta^{-},a)
   = \tau^{2}\,
     \mathrm{KL}\!\Big(
       \sigma\big(f_{\theta^{-}}(x)_{:a}/\tau\big)
       \,\Big\|\,
       \sigma\big(f_{\theta}(x)_{:a}/\tau\big)
     \Big),
\end{equation}
with $\sigma$ the softmax. Distillation is taken over \emph{replayed} samples
only, which is why the implementation refuses the configuration
$\texttt{use\_kd}\wedge\neg\texttt{use\_replay}$ outright rather than silently
distilling on the current task.

The convex weight is $\gamma_t = 1/(t+1)$, which weights every task seen so far
equally. Three alternatives (\texttt{floor}, \texttt{sqrt}, \texttt{fixed}) are
implemented and default to off; they exist because $1/(t+1)$ is principled only
when tasks are interchangeable, and under a task~0 far larger than the
increments the newest families receive a small share of the gradient. We note
also that $\gamma_t$ is meaningful only \emph{as} a convex weight against the KD
term: in an ablation that removes distillation, retaining $\gamma_t$ degenerates
into a bare $1/(t+1)$ scalar on the whole objective --- a decaying learning rate
in the costume of a replay weight. It is therefore disabled by default whenever
KD is disabled.

\paragraph{Synaptic Intelligence.}
Parameter importance is the path integral of Zenke et al. During training, each
optimiser step contributes
\begin{equation}
  \label{eq:si-accum}
  W_k \;\mathrel{+}=\; -\,g_k \cdot \Delta\theta_k,
\end{equation}
with $g_k$ the pre-step gradient and $\Delta\theta_k$ the realised parameter
change. At a task boundary the running sum is folded into the importances,
normalised by the squared distance the parameter actually travelled,
\begin{equation}
  \label{eq:si-omega}
  \Omega_k \;\mathrel{+}=\;
    \frac{W_k}{\big(\theta_k - \theta^{\text{old}}_k\big)^{2} + \xi},
  \qquad W_k \leftarrow 0,
\end{equation}
and the quadratic penalty charged during subsequent training is
\begin{equation}
  \label{eq:si-penalty}
  \Omega(\theta) \;=\; \sum_k \Omega_k \big(\theta_k - \theta^{\text{old}}_k\big)^{2}.
\end{equation}
Whether the anchor $\theta^{\text{old}}$ advances at the boundary is a separate
decision from whether $\Omega$ is updated, and Section~\ref{sec:method:order}
turns on exactly that separation.

\paragraph{Frozen normalisation statistics.}
BatchNorm layers are held in evaluation mode for the entire continual phase.
Their running statistics were estimated on task~0, and allowing a stream
dominated by the newest families to overwrite them is a channel of forgetting
that is unrelated to the weights and would otherwise be silently attributed to
the training rule.

\subsection{Diversity-aware replay buffer}
\label{sec:method:buffer}

The buffer holds at most $M$ samples and is budgeted uniformly across the
families seen so far: with $F$ distinct families held or arriving, the
per-family budget is $b = \max(1, \lfloor M / F \rfloor)$, existing families are
truncated to $b$, and each arriving family is filled to $b$.

Within a family, selection is MADAR's diversity criterion. An Isolation Forest
is fitted over the family's samples and they are ordered ascending by
$\mathrm{decision\_function}$, i.e. most anomalous first:
\begin{equation}
  \label{eq:buffer-order}
  \pi = \operatorname{argsort}\big(\mathrm{IF}(\,\cdot\,)\big),
  \qquad
  \mathcal{M}^{(f)} = \operatorname{interleave}\big(
     \pi_{1:\lfloor b/2 \rfloor},\;
     \pi_{n - \lceil b/2 \rceil + 1 : n}\big),
\end{equation}
keeping the most anomalous half and the most prototypical half. The two ends are
\emph{interleaved} rather than concatenated so that a later truncation to a
smaller budget --- which takes a prefix of the stored list --- remains balanced
between the two ends instead of retaining whichever end happens to sit first.

Two spaces are selectable for the forest. Fitting on the scaled \emph{input}
features (\texttt{--buffer-space raw}) is MADAR as published: selection is then a
fixed property of the data. Fitting on the classifier's latent representation
$h_\theta$ (\texttt{latent}, the default here; MADAR-$\theta$) makes selection
move as the model learns. The distinction is not cosmetic --- it decides whether
the memory is chosen in a fixed space or in one the learner is simultaneously
reshaping --- and it is recorded per run.\footnote{The Isolation Forest's
\texttt{contamination} parameter appears in the configuration but only shifts the
estimator's decision \emph{threshold}. Selection here is purely rank-based on
$\mathrm{decision\_function}$ scores, so its value does not change which samples
are chosen; it is retained because the estimator requires it.}

\subsection{MADAR + Unlearning}
\label{sec:method:unlearn}

The proposed condition is the backbone of Section~\ref{sec:method:madar} with a
\emph{forget phase} appended to every task $t \ge 1$. A fraction of the task the
model has just learned is nominated for deletion, the model is trained to become
uninformative about it, and the replay buffer is rebuilt so that the deleted
samples cannot return. The phase runs \emph{after} the continual phase, so the
setting is learn-then-forget: the samples in question have already been fit.
Algorithm~\ref{alg:madaru} states the whole procedure; the subsections that
follow justify the parts that are not mechanical.

\begin{algorithm}[t]
\caption{MADAR${+}$Unlearning, one task $t \ge 1$}
\label{alg:madaru}
\begin{algorithmic}[1]
  \Require model $\theta$, teacher $\theta^{-}$, buffer $\mathcal{M}_{t-1}$,
           SI state $(W,\Omega,\theta^{\text{old}})$
  \State $\theta \gets \textsc{ContinualPhase}(\theta,\theta^{-},\mathcal{D}_t,\mathcal{M}_{t-1})$
         \Comment{Eq.~\eqref{eq:cl}, $N_{\mathrm{cl}}$ steps}
  \State $\Omega \gets \textsc{FoldImportances}(W,\theta,\theta^{\text{old}})$
         \Comment{Eq.~\eqref{eq:si-omega}; anchor \emph{not} advanced}
  \State $(\mathcal{F}_t, \mathcal{R}_t) \gets \textsc{Select}(\mathcal{D}_t, \rho)$
         \Comment{Section~\ref{sec:method:select}}
  \If{$\mathcal{F}_t = \emptyset$}
      \State advance anchor, refresh teacher, rebuild buffer from $\mathcal{D}_t$;
             \textbf{return}
  \EndIf
  \State $\theta^{\text{pre}} \gets \theta$ \Comment{frozen; teacher and reference}
  \For{$e = 1$ \textbf{to} $E_u$}
    \For{minibatch $x_f \subset \mathcal{F}_t$}
      \State draw $x_r \sim \mathcal{R}_t$, \; $\tilde{x} \sim \mathcal{M}_{t-1}$
      \State $\theta \gets \theta - \eta_u \nabla_\theta \mathcal{L}^{U}$
             \Comment{Eq.~\eqref{eq:unlearn}}
    \EndFor
  \EndFor
  \State $\textsc{Measure}(\theta^{\text{pre}}, \theta, \mathcal{F}_t, \mathcal{R}_t)$
         \Comment{Section~\ref{sec:method:efficacy}}
  \State $\theta^{\text{old}} \gets \theta$;\quad $\theta^{-} \gets \theta$
  \State $\mathcal{M}_t \gets \textsc{RebuildBuffer}(\mathcal{M}_{t-1}, \mathcal{R}_t)$
         \Comment{retain split only}
\end{algorithmic}
\end{algorithm}

\subsubsection{Ordering of the SI boundary}
\label{sec:method:order}

The task boundary is closed \emph{before} the forget phase, but the anchor is
deliberately held back: Eq.~\eqref{eq:si-omega} is applied with
$\theta^{\text{old}}$ left pointing at the parameters as they stood at the
\emph{start} of task $t$, and it advances only once unlearning is complete
(lines 2 and 13 of Algorithm~\ref{alg:madaru}).

This ordering is load-bearing in both directions. Folding $W$ into $\Omega$
first means the importances already account for what task $t$ made important, so
the forget phase is constrained by a current estimate rather than a stale one.
Withholding the anchor means the penalty in Eq.~\eqref{eq:si-penalty} is
evaluated against the start of the task, so unlearning is charged for every
important parameter it drags away from that point. Had the anchor advanced
first, $\theta - \theta^{\text{old}} = 0$ at the moment the forget phase begins
and the regulariser would be identically inert for precisely the phase it is
meant to regulate.

We note for completeness that the path integral is \emph{not} accumulated during
the forget phase --- Eq.~\eqref{eq:si-accum} is applied only inside the
continual phase, since $W$ has already been consumed at line~2. The unlearning
trajectory is therefore regularised by SI but does not itself contribute to
future importance estimates.

\subsubsection{Forget-set selection}
\label{sec:method:select}

Selection partitions the current task's samples into
$\mathcal{D}_t = \mathcal{F}_t \sqcup \mathcal{R}_t$ under a budget
$k = \lfloor \rho\,|\mathcal{D}_t| \rfloor$ that is held \emph{fixed across
selectors}. Holding the budget fixed is what makes a comparison between rules
attributable to \emph{which} samples are named rather than to how many. The
seed for the decision is derived, $\mathrm{seed}_t = \mathrm{child}(\mathrm{seed}, t, 7)$,
rather than drawn from the global stream, so that introducing a selector cannot
shift every later random draw in the run and thereby de-pair two conditions
intended to share a seed. If $k = 0$ the phase is skipped and the run record
says so explicitly.

Three rules are compared.

\begin{description}
  \item[\texttt{random}] A uniform subset of size $k$. This is the control for
    \emph{which} samples: it isolates the effect of the selection rule from the
    effect of deleting $k$ samples at all.

  \item[\texttt{donut}] (default) An Isolation Forest is fitted jointly over the
    task in latent space, samples are ordered ascending by
    $\mathrm{decision\_function}$, and the forget set is the \emph{middle band}
    of that ordering,
    \begin{equation}
      \label{eq:donut}
      \mathcal{F}_t = \pi_{\,j:\,j+k},
      \qquad j = \Big\lfloor \tfrac{n}{2} \Big\rfloor - \Big\lfloor \tfrac{k}{2} \Big\rfloor .
    \end{equation}
    Both the most anomalous and the most prototypical samples are kept. This
    applies MADAR's own prior --- the criterion it uses to \emph{choose} what to
    remember, Eq.~\eqref{eq:buffer-order} --- to the complementary question of
    what to delete, and is the hand-crafted rule the learned policy is later
    measured against.

  \item[\texttt{leftover}] A per-family Isolation Forest in \emph{raw} feature
    space; each family keeps the anomalous/prototypical halves that a
    MADAR-style selection at the current per-family budget would have kept, and
    everything else is forgotten. This rule ignores $\rho$ by construction: the
    forget set is whatever the buffer budget leaves over, which is typically far
    larger than a $10\%$ slab. Comparisons involving it are therefore not
    budget-matched, and we report it as such.
\end{description}

\subsubsection{The forget objective}
\label{sec:method:objective}

Let $\theta^{\text{pre}}$ denote the model at the end of the continual phase,
frozen and used both as the distillation teacher for this phase and as the
reference for efficacy measurement. The retain anchor is therefore ``what the
model believed a moment ago'', not the previous task's model. Writing
$u = \mathbf{1}/s$ for the uniform distribution over the $s = a_t - a_{t-1}$
classes this task introduced, the phase minimises
\begin{align}
  \label{eq:unlearn}
  \mathcal{L}^{U}
   &= \alpha\,
      \underbrace{\mathrm{KL}\!\Big(u \,\Big\|\,
        \sigma\big(f_\theta(x_f)_{a_{t-1}:a_t}\big)\Big)}_{\text{forget}}
      \nonumber\\
   &\quad + (1-\alpha)\,
      \underbrace{\Big[
        \tfrac{1}{2}\,\ell_{\mathrm{CE}}\big(f_\theta(x_\cup) + m^{(a_t)},\, y_\cup\big)
        + \tfrac{1}{2}\,\mathcal{L}_{\mathrm{KD}}\big(x_\cup; \theta, \theta^{\text{pre}}, a_t\big)
      \Big]}_{\text{retain}}
      \nonumber\\
   &\quad + c\,\Omega(\theta),
\end{align}
where $x_f \subset \mathcal{F}_t$ and
$(x_\cup, y_\cup) = (x_r,y_r) \cup (\tilde{x},\tilde{y})$ concatenates a batch of
this task's \emph{retained} samples with a replay batch from the buffer. The
forget set defines the step count; the retain and replay streams cycle to match
it. Optimisation uses Adam at $\eta_u$ for $E_u$ passes over $\mathcal{F}_t$,
with gradient-norm clipping at $0.5$ and BatchNorm still frozen.

Three properties of Eq.~\eqref{eq:unlearn} should be read carefully.

\emph{The forget term acts on the new-task slice only.} The softmax is taken over
$[a_{t-1}, a_t)$ --- the $s$ classes this task introduced --- not over the full
active prefix. The objective therefore drives the model toward indifference
\emph{among the new families}, which is a narrower claim than indifference among
all classes seen so far.

\emph{Retention is anchored on two populations.} Because $x_\cup$ mixes retained
current-task samples with replayed old-family samples, the $(1-\alpha)$ term
defends both plasticity on the task being partially forgotten and stability on
everything before it, under one weight.

\emph{The regulariser is not slack.} By Section~\ref{sec:method:order}, $\Omega$
is evaluated against the start of the task, so the forget term must overcome an
active quadratic penalty on exactly those parameters the task made important.

\subsubsection{Buffer reconstruction from the retain split}
\label{sec:method:rebuild}

After the forget phase the buffer is rebuilt from $\mathcal{R}_t$ rather than
$\mathcal{D}_t$: latents are recomputed with the \emph{post}-unlearning model and
Eq.~\eqref{eq:buffer-order} is re-applied per family. Forgotten samples are
consequently never eligible for replay at any later task.

This is worth separating from the objective, because it is a second and
independent deletion mechanism, and it operates whenever a forget set is selected
--- including when $\alpha = 0$. It is the reason the ablation of
Section~\ref{sec:method:controls} is necessary rather than decorative.

\subsubsection{Efficacy measurement}
\label{sec:method:efficacy}

Unlearning is measured, per task, by comparing $\theta^{\text{pre}}$ against the
post-phase $\theta$ on the same inputs, over two disjoint scopes: the forget set
$\mathcal{F}_t$ and the retained current-task samples $\mathcal{R}_t$. For each
scope we record accuracy before and after, the mean predictive entropy and its
normalised form, and the displacement of the latent representation:
\begin{align}
  \label{eq:efficacy}
  H(x) &= -\!\!\sum_{j < a_t} p_j \log p_j,
  \qquad p = \sigma\big(f_{\theta}(x)_{:a_t}\big), \\
  \widehat{H} &= \frac{\mathbb{E}[H(x)]}{\log a_t},
  \qquad
  \Delta_{\text{lat}}(x) = \big\| h_{\theta}(x) - h_{\theta^{\text{pre}}}(x) \big\|_2 .
\end{align}
A useful result is a large drop in accuracy and a large $\Delta_{\text{lat}}$ on
$\mathcal{F}_t$ together with small movement on $\mathcal{R}_t$: forgetting that
also flattens the retained samples is collateral damage, not deletion. Note that
$\widehat{H}$ is normalised over the \emph{full} active prefix while the forget
term of Eq.~\eqref{eq:unlearn} operates on the $s$-class slice, so a forget set
driven to perfect uniformity within that slice will not register
$\widehat{H} \approx 1$; the two widths are not interchangeable and we report
$\widehat{H}$ as a monotone diagnostic rather than as a saturating one.

\subsection{Conditions, controls and ablations}
\label{sec:method:controls}

Four conditions share the network, the family selection, the scaler, the task
schedule, the masked evaluation and the determinism block. Only the rule applied
to a task differs.

\begin{table}[t]
  \centering
  \small
  \caption{The four conditions. All share every non-training component.}
  \label{tab:conditions}
  \begin{tabular}{@{}lp{0.62\linewidth}@{}}
    \toprule
    Condition & Training rule \\
    \midrule
    Naive & Fine-tune on the current task only. Lower bound. \\
    Joint & Discard the model and retrain from scratch on all data seen so far.
            Oracle upper bound; not a realisable continual setup. \\
    MADAR & Eq.~\eqref{eq:cl}: diversity-aware replay, distillation against the
            previous task's model, SI penalty. \\
    MADAR${+}$U & MADAR, plus the per-task forget phase of
            Section~\ref{sec:method:unlearn}. \\
    \bottomrule
  \end{tabular}
\end{table}

\paragraph{$\alpha = 0$ is the matched control.}
The headline comparison is not MADAR${+}$U against MADAR. At $\alpha = 0$ the
forget set is still selected by the same rule, is still excluded from the buffer
rebuild, and the same number of gradient steps is still taken --- only the push
toward uniform is absent. The contrast $\alpha > 0$ versus $\alpha = 0$ therefore
isolates the value of the \emph{forget objective}, while $\alpha = 0$ versus
MADAR isolates the value of \emph{curation}, i.e. of removing those samples from
memory at all. Reporting the method without the $\alpha = 0$ arm would attribute
to unlearning what may simply be memory curation, and both arms are run for every
corpus and seed.

\paragraph{Component ablations.}
Replay, distillation and the SI penalty are independently switchable in
Eq.~\eqref{eq:cl} so that an ablation removes exactly one thing. SI is ablated by
setting $c = 0$, which leaves the optimiser, the iteration count and the data
order untouched; removing distillation additionally disables $\gamma_t$ for the
reason given in Section~\ref{sec:method:madar}.

\subsection{Evaluation protocol}
\label{sec:method:eval}

After every task, predictions are taken as
$\hat{y} = \arg\max_{j < a_t} f_\theta(x)_j$ over the test split, and all
reported quantities are derived from a single confusion matrix so that the
numbers in a record are mutually consistent by construction rather than by three
code paths agreeing. Three scopes are evaluated:
\begin{itemize}
  \item \textbf{seen} --- all families encountered so far. The headline number.
  \item \textbf{recent} --- the families just introduced. Acquisition, i.e.
        plasticity.
  \item \textbf{task0} --- the base families. Retention, i.e. resistance to
        forgetting.
\end{itemize}
The \textbf{seen} scope averages the other two, so a method that trades
acquisition against retention is invisible in it alone; we report all three.

Macro averages are taken over classes \emph{with support in the evaluation
slice}. A class absent from a slice has no recall to measure, and scoring it as
zero would make the metric a function of which classes happen to be absent.
Because a class may be predicted without being present, the false positives such
predictions generate belong to no scored class; they are surfaced explicitly as
\texttt{n\_pred\_outside\_support} rather than dropped. For single-label
multiclass, micro-precision, micro-recall and micro-F1 all equal accuracy, so
only accuracy is reported under that name.

Two protocol constraints are enforced rather than documented. First, the
\texttt{stage1} protocol (task~0 retrained per seed) and the \texttt{meta}
protocol (a shared task-0 checkpoint) are never differenced against each other,
and are tabulated as separate blocks. Second, a crashed run is recorded as failed
and is never scored, defaulted or interpolated --- a missing number that silently
becomes \texttt{NaN} and then sorts to the top of a ranking is a real failure
mode, and the sweep boundary is where it is refused.

\subsection{Hyperparameters}
\label{sec:method:hparams}

Table~\ref{tab:hparams} gives the values in effect. Per-corpus entries are
protocol for that corpus; an explicit override always wins and, whatever wins,
the resolved value is written into the run record --- there is no value in effect
that the log does not name.

\begin{table}[t]
  \centering
  \small
  \caption{Protocol parameters. Shared values are common to every corpus.}
  \label{tab:hparams}
  \begin{tabular}{@{}llcccc@{}}
    \toprule
    & & EMBER 2018 & EMBER 2024 & LAMDA & Tiny ImageNet \\
    \midrule
    \multicolumn{6}{@{}l}{\emph{Shared}}\\
    Task-0 epochs $E_0$        & \texttt{task0\_epochs}   & \multicolumn{3}{c}{30, Adam $10^{-3}$} & 50 \\
    Batch size                 & \texttt{batch\_size}     & \multicolumn{3}{c}{256} & 32 \\
    Continual optimiser        & ---                      & \multicolumn{4}{c}{SGD, momentum $0.9$, weight decay $10^{-6}$, clip $1.0$} \\
    Distillation temperature   & \texttt{kd\_temp}        & \multicolumn{4}{c}{$\tau = 2.0$} \\
    Current-task weight        & \texttt{rnt\_mode}       & \multicolumn{4}{c}{$\gamma_t = 1/(t{+}1)$} \\
    Forget weight              & \texttt{alpha}           & \multicolumn{4}{c}{$\alpha \in \{0.2,\ 0\}$ (both arms)} \\
    Forget budget              & \texttt{forget\_ratio}   & \multicolumn{4}{c}{$\rho = 0.10$} \\
    Forget passes              & \texttt{unlearn\_epochs} & \multicolumn{3}{c}{$E_u = 3$} & 3 \\
    Forget optimiser           & \texttt{unlearn\_lr}     & \multicolumn{4}{c}{Adam $\eta_u = 10^{-4}$, clip $0.5$} \\
    Family eligibility         & \texttt{min\_family\_samples} & \multicolumn{3}{c}{$\ge 200$} & --- \\
    \midrule
    \multicolumn{6}{@{}l}{\emph{Per corpus}}\\
    Schedule                   & \texttt{--task-setup}    & $50{+}5{\times}10$ & $50{+}5{\times}10$ & $30{+}5{\times}10$ & $100{+}10{\times}10$ \\
    Continual budget $N_{\mathrm{cl}}$ & \texttt{cl\_iters} & 2000 & 8000 & 4000 & 3000 \\
    Continual lr               & \texttt{cl\_lr}          & $10^{-4}$ & $10^{-4}$ & $10^{-4}$ & $0.03$ \\
    Buffer capacity $M$        & \texttt{mem\_size}       & 5000 & 10000 & 2500 & 2000 \\
    SI strength $c$            & \texttt{si\_c}           & 1.0 & 100.0 & 100.0 & 1.0 \\
    Per-family cap             & \texttt{family\_cap}     & --- & 10000 & --- & --- \\
    Feature clip               & \texttt{feature\_clip}   & $\pm10$ & $\pm10$ & --- (binary) & --- \\
    \bottomrule
  \end{tabular}
\end{table}

\subsection{Threats to validity}
\label{sec:method:threats}

\paragraph{Compute is not matched between MADAR and MADAR${+}$U.}
The forget phase adds $E_u \lceil |\mathcal{F}_t| / B \rceil$ optimiser steps per
task on top of the $N_{\mathrm{cl}}$ continual steps, and the run record reports
the total. The difference is small relative to $N_{\mathrm{cl}}$ but it is not
zero, so a MADAR${+}$U improvement over MADAR is confounded with additional
gradient steps. The $\alpha = 0$ arm of Section~\ref{sec:method:controls} is
matched in step count and is the comparison from which the forget objective's
contribution should be read.

\paragraph{The forget objective and the efficacy metric span different widths.}
As noted in Section~\ref{sec:method:efficacy}, $\mathcal{L}^U$ flattens the
$s$-class new-task slice while $\widehat{H}$ normalises over the full active
prefix of width $a_t$. The two are consistent in direction but not on a common
scale.

\paragraph{Unlearning is behavioural, not certified.}
Eq.~\eqref{eq:unlearn} drives the model's \emph{predictions} on $\mathcal{F}_t$
toward uniformity and removes those samples from memory; it offers no guarantee
of indistinguishability from a model retrained without them, and we make no such
claim. The reported quantities of Section~\ref{sec:method:efficacy} are evidence
of behavioural forgetting under a specific probe, and the \texttt{Joint} oracle
of Table~\ref{tab:conditions} is a retraining reference for accuracy, not a
privacy reference.

\paragraph{The \texttt{leftover} selector is not budget-matched.}
Its forget set is determined by the buffer budget rather than by $\rho$, so
differences against \texttt{random} and \texttt{donut} confound the selection
rule with the number of samples deleted.

\paragraph{Reproduction.}
Every reported run is produced by a single command and writes one self-describing
JSON record naming the corpus, the schedule, the model and parameter count, every
hyperparameter in effect, the seed, the protocol, the library versions, the GPU
and the git commit:
\begin{verbatim}
./iclr run --dataset ember2018 --task-setup ember \
           --experiment madar_unlearn --selector donut \
           --model full --seed 1 --group <group>
\end{verbatim}

% or compile the preview wrapper docs/methodology_standalone.tex.
%
% Required packages (see the wrapper for a working preamble):
%   amsmath, amssymb, booktabs, algorithm, algpseudocode, xcolor, hyperref
%
% Every equation, constant and default below is transcribed from the
% implementation in code/; the trailing comment on a line names the source.
% ---------------------------------------------------------------------------

\section{Methodology}
\label{sec:method}

\subsection{Problem setting and notation}
\label{sec:method:notation}

We study class-incremental learning over a stream of $T{+}1$ tasks. A task
setup is written $C_0 + s \times T$: task $0$ introduces $C_0$ base families and
each subsequent task $t \in \{1,\dots,T\}$ introduces $s$ previously unseen
families, for $K = C_0 + sT$ classes in total.\footnote{\texttt{core/tasks.py}.
The EMBER protocol is $50+5\times10$ ($K=100$); LAMDA is $30+5\times10$
($K=80$).} Class indices follow \emph{presentation order}, so the classes
introduced by task $t$ are
\begin{equation}
  \mathcal{C}_t =
  \begin{cases}
    \{0,\dots,C_0-1\}, & t = 0,\\[2pt]
    \{C_0 + (t-1)s, \dots, C_0 + ts - 1\}, & t \ge 1,
  \end{cases}
\end{equation}
and the \emph{active prefix} after task $t$ has width $a_t = C_0 + ts$, with
$a_{-1} \coloneqq C_0$ by convention. Ordering the classes this way is what makes
prefix masking well defined: it is not a cosmetic relabelling.

Let $f_\theta : \mathcal{X} \to \mathbb{R}^{K}$ be the classifier, emitting a
logit vector of full width $K$ at every stage, and let
$h_\theta : \mathcal{X} \to \mathbb{R}^{d_{\text{lat}}}$ be its penultimate
(latent) representation. Classes that have not yet arrived are suppressed by an
additive mask rather than by reshaping the head,
\begin{equation}
  \label{eq:mask}
  m^{(a)} \in \mathbb{R}^{K}, \qquad
  m^{(a)}_j = \begin{cases} 0, & j < a, \\ -10^{9}, & j \ge a, \end{cases}
\end{equation}
so a single parameterisation carries the whole stream and no logit belonging to
an unseen class can contribute to a loss or to an $\arg\max$.%
\footnote{\texttt{core/training.py:logit\_mask}.}
Table~\ref{tab:notation} collects the remaining symbols.

\begin{table}[t]
  \centering
  \small
  \caption{Notation, and where each quantity is realised in the implementation.}
  \label{tab:notation}
  \begin{tabular}{@{}llp{0.30\linewidth}@{}}
    \toprule
    Symbol & Meaning & Implementation \\
    \midrule
    $\mathcal{D}_t$ & training samples of task $t$ & \texttt{Experiment.task\_loader} \\
    $a_t$ & active logit width after task $t$ & \texttt{TaskSchedule.active\_count} \\
    $m^{(a)}$ & additive prefix mask, Eq.~\eqref{eq:mask} & \texttt{training.logit\_mask} \\
    $\mathcal{M}_t$ & replay buffer after task $t$ & \texttt{core.buffer.ReplayBuffer} \\
    $M$ & total buffer capacity & \texttt{mem\_size} \\
    $\theta^{-}$ & frozen teacher & \texttt{training.clone\_teacher} \\
    $\theta^{\text{pre}}_t$ & model before the forget phase & \texttt{pre} in \texttt{train\_task} \\
    $\mathcal{F}_t,\ \mathcal{R}_t$ & forget / retain split of $\mathcal{D}_t$ & \texttt{select\_forget\_set} \\
    $\rho$ & forget budget as a fraction of $|\mathcal{D}_t|$ & \texttt{forget\_ratio} \\
    $\alpha$ & weight on the forget term & \texttt{alpha} \\
    $\gamma_t$ & weight on the current-task term & \texttt{training.rnt\_for} \\
    $\Omega,\ \theta^{\text{old}}$ & SI importances and anchor & \texttt{SynapticIntelligence} \\
    $c$ & SI strength & \texttt{si\_c} \\
    $\tau$ & distillation temperature & \texttt{kd\_temp} \\
    \bottomrule
  \end{tabular}
\end{table}

\paragraph{Preprocessing.}
Families enter the study only if they carry at least $200$ training samples, and
the $K$ most frequent eligible families are retained. Feature standardisation is
fit \emph{on task~0 alone} and applied to the whole stream, clipped to
$\pm 10$ standard deviations; fitting on the full training set would leak the
statistics of families the learner has not yet met into a protocol whose entire
claim is that it has not met them. LAMDA's features are binary and are left
unscaled.

\subsection{Shared continual backbone}
\label{sec:method:madar}

All conditions share one training loop, one evaluation protocol and one
determinism block, and differ only in the rule applied to a task
(\texttt{experiments/base.py}). We describe that shared backbone first, because
MADAR${+}$Unlearning is defined as a strict extension of it: everything in this
subsection is common to both, so any difference measured between them is
attributable to Section~\ref{sec:method:unlearn} alone.

\paragraph{Task 0.}
The base families are fit by ordinary supervised training for $E_0$ epochs with
Adam,
\begin{equation}
  \label{eq:task0}
  \min_\theta\ \mathbb{E}_{(x,y)\sim\mathcal{D}_0}
    \big[\ \ell_{\mathrm{CE}}\!\big(f_\theta(x) + m^{(a_0)},\, y\big)\ \big].
\end{equation}
Afterwards the SI anchor is set, $\theta^{\text{old}} \leftarrow \theta$, the
teacher is cloned, $\theta^{-} \leftarrow \theta$, and the buffer is filled from
$\mathcal{D}_0$ by the rule of Section~\ref{sec:method:buffer}.

\paragraph{Continual phase.}
For each $t \ge 1$ the model is trained for a \emph{fixed iteration budget}
$N_{\mathrm{cl}}$ rather than a fixed number of epochs, so that compute per task
does not track task size. Each step draws one batch $(x,y) \sim \mathcal{D}_t$
and one replay batch $(\tilde{x},\tilde{y}) \sim \mathcal{M}_{t-1}$ and minimises
\begin{align}
  \label{eq:cl}
  \mathcal{L}^{\mathrm{CL}}_t
   &= \gamma_t\,
      \ell_{\mathrm{CE}}\!\big(f_\theta(x \cup \tilde{x}) + m^{(a_t)},\;
                                y \cup \tilde{y}\big) \nonumber\\
   &\quad + (1-\gamma_t)\,
      \mathcal{L}_{\mathrm{KD}}\!\big(\tilde{x};\, \theta,\, \theta^{-},\, a_{t-1}\big)
      \;+\; c\,\Omega(\theta),
\end{align}
where the distillation term is the temperature-scaled divergence \emph{from} the
teacher \emph{to} the student, restricted to the previous active prefix,
\begin{equation}
  \label{eq:kd}
  \mathcal{L}_{\mathrm{KD}}(x;\theta,\theta^{-},a)
   = \tau^{2}\,
     \mathrm{KL}\!\Big(
       \sigma\big(f_{\theta^{-}}(x)_{:a}/\tau\big)
       \,\Big\|\,
       \sigma\big(f_{\theta}(x)_{:a}/\tau\big)
     \Big),
\end{equation}
with $\sigma$ the softmax. Distillation is taken over \emph{replayed} samples
only, which is why the implementation refuses the configuration
$\texttt{use\_kd}\wedge\neg\texttt{use\_replay}$ outright rather than silently
distilling on the current task.

The convex weight is $\gamma_t = 1/(t+1)$, which weights every task seen so far
equally. Three alternatives (\texttt{floor}, \texttt{sqrt}, \texttt{fixed}) are
implemented and default to off; they exist because $1/(t+1)$ is principled only
when tasks are interchangeable, and under a task~0 far larger than the
increments the newest families receive a small share of the gradient. We note
also that $\gamma_t$ is meaningful only \emph{as} a convex weight against the KD
term: in an ablation that removes distillation, retaining $\gamma_t$ degenerates
into a bare $1/(t+1)$ scalar on the whole objective --- a decaying learning rate
in the costume of a replay weight. It is therefore disabled by default whenever
KD is disabled.

\paragraph{Synaptic Intelligence.}
Parameter importance is the path integral of Zenke et al. During training, each
optimiser step contributes
\begin{equation}
  \label{eq:si-accum}
  W_k \;\mathrel{+}=\; -\,g_k \cdot \Delta\theta_k,
\end{equation}
with $g_k$ the pre-step gradient and $\Delta\theta_k$ the realised parameter
change. At a task boundary the running sum is folded into the importances,
normalised by the squared distance the parameter actually travelled,
\begin{equation}
  \label{eq:si-omega}
  \Omega_k \;\mathrel{+}=\;
    \frac{W_k}{\big(\theta_k - \theta^{\text{old}}_k\big)^{2} + \xi},
  \qquad W_k \leftarrow 0,
\end{equation}
and the quadratic penalty charged during subsequent training is
\begin{equation}
  \label{eq:si-penalty}
  \Omega(\theta) \;=\; \sum_k \Omega_k \big(\theta_k - \theta^{\text{old}}_k\big)^{2}.
\end{equation}
Whether the anchor $\theta^{\text{old}}$ advances at the boundary is a separate
decision from whether $\Omega$ is updated, and Section~\ref{sec:method:order}
turns on exactly that separation.

\paragraph{Frozen normalisation statistics.}
BatchNorm layers are held in evaluation mode for the entire continual phase.
Their running statistics were estimated on task~0, and allowing a stream
dominated by the newest families to overwrite them is a channel of forgetting
that is unrelated to the weights and would otherwise be silently attributed to
the training rule.

\subsection{Diversity-aware replay buffer}
\label{sec:method:buffer}

The buffer holds at most $M$ samples and is budgeted uniformly across the
families seen so far: with $F$ distinct families held or arriving, the
per-family budget is $b = \max(1, \lfloor M / F \rfloor)$, existing families are
truncated to $b$, and each arriving family is filled to $b$.

Within a family, selection is MADAR's diversity criterion. An Isolation Forest
is fitted over the family's samples and they are ordered ascending by
$\mathrm{decision\_function}$, i.e. most anomalous first:
\begin{equation}
  \label{eq:buffer-order}
  \pi = \operatorname{argsort}\big(\mathrm{IF}(\,\cdot\,)\big),
  \qquad
  \mathcal{M}^{(f)} = \operatorname{interleave}\big(
     \pi_{1:\lfloor b/2 \rfloor},\;
     \pi_{n - \lceil b/2 \rceil + 1 : n}\big),
\end{equation}
keeping the most anomalous half and the most prototypical half. The two ends are
\emph{interleaved} rather than concatenated so that a later truncation to a
smaller budget --- which takes a prefix of the stored list --- remains balanced
between the two ends instead of retaining whichever end happens to sit first.

Two spaces are selectable for the forest. Fitting on the scaled \emph{input}
features (\texttt{--buffer-space raw}) is MADAR as published: selection is then a
fixed property of the data. Fitting on the classifier's latent representation
$h_\theta$ (\texttt{latent}, the default here; MADAR-$\theta$) makes selection
move as the model learns. The distinction is not cosmetic --- it decides whether
the memory is chosen in a fixed space or in one the learner is simultaneously
reshaping --- and it is recorded per run.\footnote{The Isolation Forest's
\texttt{contamination} parameter appears in the configuration but only shifts the
estimator's decision \emph{threshold}. Selection here is purely rank-based on
$\mathrm{decision\_function}$ scores, so its value does not change which samples
are chosen; it is retained because the estimator requires it.}

\subsection{MADAR + Unlearning}
\label{sec:method:unlearn}

The proposed condition is the backbone of Section~\ref{sec:method:madar} with a
\emph{forget phase} appended to every task $t \ge 1$. A fraction of the task the
model has just learned is nominated for deletion, the model is trained to become
uninformative about it, and the replay buffer is rebuilt so that the deleted
samples cannot return. The phase runs \emph{after} the continual phase, so the
setting is learn-then-forget: the samples in question have already been fit.
Algorithm~\ref{alg:madaru} states the whole procedure; the subsections that
follow justify the parts that are not mechanical.

\begin{algorithm}[t]
\caption{MADAR${+}$Unlearning, one task $t \ge 1$}
\label{alg:madaru}
\begin{algorithmic}[1]
  \Require model $\theta$, teacher $\theta^{-}$, buffer $\mathcal{M}_{t-1}$,
           SI state $(W,\Omega,\theta^{\text{old}})$
  \State $\theta \gets \textsc{ContinualPhase}(\theta,\theta^{-},\mathcal{D}_t,\mathcal{M}_{t-1})$
         \Comment{Eq.~\eqref{eq:cl}, $N_{\mathrm{cl}}$ steps}
  \State $\Omega \gets \textsc{FoldImportances}(W,\theta,\theta^{\text{old}})$
         \Comment{Eq.~\eqref{eq:si-omega}; anchor \emph{not} advanced}
  \State $(\mathcal{F}_t, \mathcal{R}_t) \gets \textsc{Select}(\mathcal{D}_t, \rho)$
         \Comment{Section~\ref{sec:method:select}}
  \If{$\mathcal{F}_t = \emptyset$}
      \State advance anchor, refresh teacher, rebuild buffer from $\mathcal{D}_t$;
             \textbf{return}
  \EndIf
  \State $\theta^{\text{pre}} \gets \theta$ \Comment{frozen; teacher and reference}
  \For{$e = 1$ \textbf{to} $E_u$}
    \For{minibatch $x_f \subset \mathcal{F}_t$}
      \State draw $x_r \sim \mathcal{R}_t$, \; $\tilde{x} \sim \mathcal{M}_{t-1}$
      \State $\theta \gets \theta - \eta_u \nabla_\theta \mathcal{L}^{U}$
             \Comment{Eq.~\eqref{eq:unlearn}}
    \EndFor
  \EndFor
  \State $\textsc{Measure}(\theta^{\text{pre}}, \theta, \mathcal{F}_t, \mathcal{R}_t)$
         \Comment{Section~\ref{sec:method:efficacy}}
  \State $\theta^{\text{old}} \gets \theta$;\quad $\theta^{-} \gets \theta$
  \State $\mathcal{M}_t \gets \textsc{RebuildBuffer}(\mathcal{M}_{t-1}, \mathcal{R}_t)$
         \Comment{retain split only}
\end{algorithmic}
\end{algorithm}

\subsubsection{Ordering of the SI boundary}
\label{sec:method:order}

The task boundary is closed \emph{before} the forget phase, but the anchor is
deliberately held back: Eq.~\eqref{eq:si-omega} is applied with
$\theta^{\text{old}}$ left pointing at the parameters as they stood at the
\emph{start} of task $t$, and it advances only once unlearning is complete
(lines 2 and 13 of Algorithm~\ref{alg:madaru}).

This ordering is load-bearing in both directions. Folding $W$ into $\Omega$
first means the importances already account for what task $t$ made important, so
the forget phase is constrained by a current estimate rather than a stale one.
Withholding the anchor means the penalty in Eq.~\eqref{eq:si-penalty} is
evaluated against the start of the task, so unlearning is charged for every
important parameter it drags away from that point. Had the anchor advanced
first, $\theta - \theta^{\text{old}} = 0$ at the moment the forget phase begins
and the regulariser would be identically inert for precisely the phase it is
meant to regulate.

We note for completeness that the path integral is \emph{not} accumulated during
the forget phase --- Eq.~\eqref{eq:si-accum} is applied only inside the
continual phase, since $W$ has already been consumed at line~2. The unlearning
trajectory is therefore regularised by SI but does not itself contribute to
future importance estimates.

\subsubsection{Forget-set selection}
\label{sec:method:select}

Selection partitions the current task's samples into
$\mathcal{D}_t = \mathcal{F}_t \sqcup \mathcal{R}_t$ under a budget
$k = \lfloor \rho\,|\mathcal{D}_t| \rfloor$ that is held \emph{fixed across
selectors}. Holding the budget fixed is what makes a comparison between rules
attributable to \emph{which} samples are named rather than to how many. The
seed for the decision is derived, $\mathrm{seed}_t = \mathrm{child}(\mathrm{seed}, t, 7)$,
rather than drawn from the global stream, so that introducing a selector cannot
shift every later random draw in the run and thereby de-pair two conditions
intended to share a seed. If $k = 0$ the phase is skipped and the run record
says so explicitly.

Three rules are compared.

\begin{description}
  \item[\texttt{random}] A uniform subset of size $k$. This is the control for
    \emph{which} samples: it isolates the effect of the selection rule from the
    effect of deleting $k$ samples at all.

  \item[\texttt{donut}] (default) An Isolation Forest is fitted jointly over the
    task in latent space, samples are ordered ascending by
    $\mathrm{decision\_function}$, and the forget set is the \emph{middle band}
    of that ordering,
    \begin{equation}
      \label{eq:donut}
      \mathcal{F}_t = \pi_{\,j:\,j+k},
      \qquad j = \Big\lfloor \tfrac{n}{2} \Big\rfloor - \Big\lfloor \tfrac{k}{2} \Big\rfloor .
    \end{equation}
    Both the most anomalous and the most prototypical samples are kept. This
    applies MADAR's own prior --- the criterion it uses to \emph{choose} what to
    remember, Eq.~\eqref{eq:buffer-order} --- to the complementary question of
    what to delete, and is the hand-crafted rule the learned policy is later
    measured against.

  \item[\texttt{leftover}] A per-family Isolation Forest in \emph{raw} feature
    space; each family keeps the anomalous/prototypical halves that a
    MADAR-style selection at the current per-family budget would have kept, and
    everything else is forgotten. This rule ignores $\rho$ by construction: the
    forget set is whatever the buffer budget leaves over, which is typically far
    larger than a $10\%$ slab. Comparisons involving it are therefore not
    budget-matched, and we report it as such.
\end{description}

\subsubsection{The forget objective}
\label{sec:method:objective}

Let $\theta^{\text{pre}}$ denote the model at the end of the continual phase,
frozen and used both as the distillation teacher for this phase and as the
reference for efficacy measurement. The retain anchor is therefore ``what the
model believed a moment ago'', not the previous task's model. Writing
$u = \mathbf{1}/s$ for the uniform distribution over the $s = a_t - a_{t-1}$
classes this task introduced, the phase minimises
\begin{align}
  \label{eq:unlearn}
  \mathcal{L}^{U}
   &= \alpha\,
      \underbrace{\mathrm{KL}\!\Big(u \,\Big\|\,
        \sigma\big(f_\theta(x_f)_{a_{t-1}:a_t}\big)\Big)}_{\text{forget}}
      \nonumber\\
   &\quad + (1-\alpha)\,
      \underbrace{\Big[
        \tfrac{1}{2}\,\ell_{\mathrm{CE}}\big(f_\theta(x_\cup) + m^{(a_t)},\, y_\cup\big)
        + \tfrac{1}{2}\,\mathcal{L}_{\mathrm{KD}}\big(x_\cup; \theta, \theta^{\text{pre}}, a_t\big)
      \Big]}_{\text{retain}}
      \nonumber\\
   &\quad + c\,\Omega(\theta),
\end{align}
where $x_f \subset \mathcal{F}_t$ and
$(x_\cup, y_\cup) = (x_r,y_r) \cup (\tilde{x},\tilde{y})$ concatenates a batch of
this task's \emph{retained} samples with a replay batch from the buffer. The
forget set defines the step count; the retain and replay streams cycle to match
it. Optimisation uses Adam at $\eta_u$ for $E_u$ passes over $\mathcal{F}_t$,
with gradient-norm clipping at $0.5$ and BatchNorm still frozen.

Three properties of Eq.~\eqref{eq:unlearn} should be read carefully.

\emph{The forget term acts on the new-task slice only.} The softmax is taken over
$[a_{t-1}, a_t)$ --- the $s$ classes this task introduced --- not over the full
active prefix. The objective therefore drives the model toward indifference
\emph{among the new families}, which is a narrower claim than indifference among
all classes seen so far.

\emph{Retention is anchored on two populations.} Because $x_\cup$ mixes retained
current-task samples with replayed old-family samples, the $(1-\alpha)$ term
defends both plasticity on the task being partially forgotten and stability on
everything before it, under one weight.

\emph{The regulariser is not slack.} By Section~\ref{sec:method:order}, $\Omega$
is evaluated against the start of the task, so the forget term must overcome an
active quadratic penalty on exactly those parameters the task made important.

\subsubsection{Buffer reconstruction from the retain split}
\label{sec:method:rebuild}

After the forget phase the buffer is rebuilt from $\mathcal{R}_t$ rather than
$\mathcal{D}_t$: latents are recomputed with the \emph{post}-unlearning model and
Eq.~\eqref{eq:buffer-order} is re-applied per family. Forgotten samples are
consequently never eligible for replay at any later task.

This is worth separating from the objective, because it is a second and
independent deletion mechanism, and it operates whenever a forget set is selected
--- including when $\alpha = 0$. It is the reason the ablation of
Section~\ref{sec:method:controls} is necessary rather than decorative.

\subsubsection{Efficacy measurement}
\label{sec:method:efficacy}

Unlearning is measured, per task, by comparing $\theta^{\text{pre}}$ against the
post-phase $\theta$ on the same inputs, over two disjoint scopes: the forget set
$\mathcal{F}_t$ and the retained current-task samples $\mathcal{R}_t$. For each
scope we record accuracy before and after, the mean predictive entropy and its
normalised form, and the displacement of the latent representation:
\begin{align}
  \label{eq:efficacy}
  H(x) &= -\!\!\sum_{j < a_t} p_j \log p_j,
  \qquad p = \sigma\big(f_{\theta}(x)_{:a_t}\big), \\
  \widehat{H} &= \frac{\mathbb{E}[H(x)]}{\log a_t},
  \qquad
  \Delta_{\text{lat}}(x) = \big\| h_{\theta}(x) - h_{\theta^{\text{pre}}}(x) \big\|_2 .
\end{align}
A useful result is a large drop in accuracy and a large $\Delta_{\text{lat}}$ on
$\mathcal{F}_t$ together with small movement on $\mathcal{R}_t$: forgetting that
also flattens the retained samples is collateral damage, not deletion. Note that
$\widehat{H}$ is normalised over the \emph{full} active prefix while the forget
term of Eq.~\eqref{eq:unlearn} operates on the $s$-class slice, so a forget set
driven to perfect uniformity within that slice will not register
$\widehat{H} \approx 1$; the two widths are not interchangeable and we report
$\widehat{H}$ as a monotone diagnostic rather than as a saturating one.

\subsection{Conditions, controls and ablations}
\label{sec:method:controls}

Four conditions share the network, the family selection, the scaler, the task
schedule, the masked evaluation and the determinism block. Only the rule applied
to a task differs.

\begin{table}[t]
  \centering
  \small
  \caption{The four conditions. All share every non-training component.}
  \label{tab:conditions}
  \begin{tabular}{@{}lp{0.62\linewidth}@{}}
    \toprule
    Condition & Training rule \\
    \midrule
    Naive & Fine-tune on the current task only. Lower bound. \\
    Joint & Discard the model and retrain from scratch on all data seen so far.
            Oracle upper bound; not a realisable continual setup. \\
    MADAR & Eq.~\eqref{eq:cl}: diversity-aware replay, distillation against the
            previous task's model, SI penalty. \\
    MADAR${+}$U & MADAR, plus the per-task forget phase of
            Section~\ref{sec:method:unlearn}. \\
    \bottomrule
  \end{tabular}
\end{table}

\paragraph{$\alpha = 0$ is the matched control.}
The headline comparison is not MADAR${+}$U against MADAR. At $\alpha = 0$ the
forget set is still selected by the same rule, is still excluded from the buffer
rebuild, and the same number of gradient steps is still taken --- only the push
toward uniform is absent. The contrast $\alpha > 0$ versus $\alpha = 0$ therefore
isolates the value of the \emph{forget objective}, while $\alpha = 0$ versus
MADAR isolates the value of \emph{curation}, i.e. of removing those samples from
memory at all. Reporting the method without the $\alpha = 0$ arm would attribute
to unlearning what may simply be memory curation, and both arms are run for every
corpus and seed.

\paragraph{Component ablations.}
Replay, distillation and the SI penalty are independently switchable in
Eq.~\eqref{eq:cl} so that an ablation removes exactly one thing. SI is ablated by
setting $c = 0$, which leaves the optimiser, the iteration count and the data
order untouched; removing distillation additionally disables $\gamma_t$ for the
reason given in Section~\ref{sec:method:madar}.

\subsection{Evaluation protocol}
\label{sec:method:eval}

After every task, predictions are taken as
$\hat{y} = \arg\max_{j < a_t} f_\theta(x)_j$ over the test split, and all
reported quantities are derived from a single confusion matrix so that the
numbers in a record are mutually consistent by construction rather than by three
code paths agreeing. Three scopes are evaluated:
\begin{itemize}
  \item \textbf{seen} --- all families encountered so far. The headline number.
  \item \textbf{recent} --- the families just introduced. Acquisition, i.e.
        plasticity.
  \item \textbf{task0} --- the base families. Retention, i.e. resistance to
        forgetting.
\end{itemize}
The \textbf{seen} scope averages the other two, so a method that trades
acquisition against retention is invisible in it alone; we report all three.

Macro averages are taken over classes \emph{with support in the evaluation
slice}. A class absent from a slice has no recall to measure, and scoring it as
zero would make the metric a function of which classes happen to be absent.
Because a class may be predicted without being present, the false positives such
predictions generate belong to no scored class; they are surfaced explicitly as
\texttt{n\_pred\_outside\_support} rather than dropped. For single-label
multiclass, micro-precision, micro-recall and micro-F1 all equal accuracy, so
only accuracy is reported under that name.

Two protocol constraints are enforced rather than documented. First, the
\texttt{stage1} protocol (task~0 retrained per seed) and the \texttt{meta}
protocol (a shared task-0 checkpoint) are never differenced against each other,
and are tabulated as separate blocks. Second, a crashed run is recorded as failed
and is never scored, defaulted or interpolated --- a missing number that silently
becomes \texttt{NaN} and then sorts to the top of a ranking is a real failure
mode, and the sweep boundary is where it is refused.

\subsection{Hyperparameters}
\label{sec:method:hparams}

Table~\ref{tab:hparams} gives the values in effect. Per-corpus entries are
protocol for that corpus; an explicit override always wins and, whatever wins,
the resolved value is written into the run record --- there is no value in effect
that the log does not name.

\begin{table}[t]
  \centering
  \small
  \caption{Protocol parameters. Shared values are common to every corpus.}
  \label{tab:hparams}
  \begin{tabular}{@{}llcccc@{}}
    \toprule
    & & EMBER 2018 & EMBER 2024 & LAMDA & Tiny ImageNet \\
    \midrule
    \multicolumn{6}{@{}l}{\emph{Shared}}\\
    Task-0 epochs $E_0$        & \texttt{task0\_epochs}   & \multicolumn{3}{c}{30, Adam $10^{-3}$} & 50 \\
    Batch size                 & \texttt{batch\_size}     & \multicolumn{3}{c}{256} & 32 \\
    Continual optimiser        & ---                      & \multicolumn{4}{c}{SGD, momentum $0.9$, weight decay $10^{-6}$, clip $1.0$} \\
    Distillation temperature   & \texttt{kd\_temp}        & \multicolumn{4}{c}{$\tau = 2.0$} \\
    Current-task weight        & \texttt{rnt\_mode}       & \multicolumn{4}{c}{$\gamma_t = 1/(t{+}1)$} \\
    Forget weight              & \texttt{alpha}           & \multicolumn{4}{c}{$\alpha \in \{0.2,\ 0\}$ (both arms)} \\
    Forget budget              & \texttt{forget\_ratio}   & \multicolumn{4}{c}{$\rho = 0.10$} \\
    Forget passes              & \texttt{unlearn\_epochs} & \multicolumn{3}{c}{$E_u = 3$} & 3 \\
    Forget optimiser           & \texttt{unlearn\_lr}     & \multicolumn{4}{c}{Adam $\eta_u = 10^{-4}$, clip $0.5$} \\
    Family eligibility         & \texttt{min\_family\_samples} & \multicolumn{3}{c}{$\ge 200$} & --- \\
    \midrule
    \multicolumn{6}{@{}l}{\emph{Per corpus}}\\
    Schedule                   & \texttt{--task-setup}    & $50{+}5{\times}10$ & $50{+}5{\times}10$ & $30{+}5{\times}10$ & $100{+}10{\times}10$ \\
    Continual budget $N_{\mathrm{cl}}$ & \texttt{cl\_iters} & 2000 & 8000 & 4000 & 3000 \\
    Continual lr               & \texttt{cl\_lr}          & $10^{-4}$ & $10^{-4}$ & $10^{-4}$ & $0.03$ \\
    Buffer capacity $M$        & \texttt{mem\_size}       & 5000 & 10000 & 2500 & 2000 \\
    SI strength $c$            & \texttt{si\_c}           & 1.0 & 100.0 & 100.0 & 1.0 \\
    Per-family cap             & \texttt{family\_cap}     & --- & 10000 & --- & --- \\
    Feature clip               & \texttt{feature\_clip}   & $\pm10$ & $\pm10$ & --- (binary) & --- \\
    \bottomrule
  \end{tabular}
\end{table}

\subsection{Threats to validity}
\label{sec:method:threats}

\paragraph{Compute is not matched between MADAR and MADAR${+}$U.}
The forget phase adds $E_u \lceil |\mathcal{F}_t| / B \rceil$ optimiser steps per
task on top of the $N_{\mathrm{cl}}$ continual steps, and the run record reports
the total. The difference is small relative to $N_{\mathrm{cl}}$ but it is not
zero, so a MADAR${+}$U improvement over MADAR is confounded with additional
gradient steps. The $\alpha = 0$ arm of Section~\ref{sec:method:controls} is
matched in step count and is the comparison from which the forget objective's
contribution should be read.

\paragraph{The forget objective and the efficacy metric span different widths.}
As noted in Section~\ref{sec:method:efficacy}, $\mathcal{L}^U$ flattens the
$s$-class new-task slice while $\widehat{H}$ normalises over the full active
prefix of width $a_t$. The two are consistent in direction but not on a common
scale.

\paragraph{Unlearning is behavioural, not certified.}
Eq.~\eqref{eq:unlearn} drives the model's \emph{predictions} on $\mathcal{F}_t$
toward uniformity and removes those samples from memory; it offers no guarantee
of indistinguishability from a model retrained without them, and we make no such
claim. The reported quantities of Section~\ref{sec:method:efficacy} are evidence
of behavioural forgetting under a specific probe, and the \texttt{Joint} oracle
of Table~\ref{tab:conditions} is a retraining reference for accuracy, not a
privacy reference.

\paragraph{The \texttt{leftover} selector is not budget-matched.}
Its forget set is determined by the buffer budget rather than by $\rho$, so
differences against \texttt{random} and \texttt{donut} confound the selection
rule with the number of samples deleted.

\paragraph{Reproduction.}
Every reported run is produced by a single command and writes one self-describing
JSON record naming the corpus, the schedule, the model and parameter count, every
hyperparameter in effect, the seed, the protocol, the library versions, the GPU
and the git commit:
\begin{verbatim}
./iclr run --dataset ember2018 --task-setup ember \
           --experiment madar_unlearn --selector donut \
           --model full --seed 1 --group <group>
\end{verbatim}

% or compile the preview wrapper docs/methodology_standalone.tex.
%
% Required packages (see the wrapper for a working preamble):
%   amsmath, amssymb, booktabs, algorithm, algpseudocode, xcolor, hyperref
%
% Every equation, constant and default below is transcribed from the
% implementation in code/; the trailing comment on a line names the source.
% ---------------------------------------------------------------------------

\section{Methodology}
\label{sec:method}

\subsection{Problem setting and notation}
\label{sec:method:notation}

We study class-incremental learning over a stream of $T{+}1$ tasks. A task
setup is written $C_0 + s \times T$: task $0$ introduces $C_0$ base families and
each subsequent task $t \in \{1,\dots,T\}$ introduces $s$ previously unseen
families, for $K = C_0 + sT$ classes in total.\footnote{\texttt{core/tasks.py}.
The EMBER protocol is $50+5\times10$ ($K=100$); LAMDA is $30+5\times10$
($K=80$).} Class indices follow \emph{presentation order}, so the classes
introduced by task $t$ are
\begin{equation}
  \mathcal{C}_t =
  \begin{cases}
    \{0,\dots,C_0-1\}, & t = 0,\\[2pt]
    \{C_0 + (t-1)s, \dots, C_0 + ts - 1\}, & t \ge 1,
  \end{cases}
\end{equation}
and the \emph{active prefix} after task $t$ has width $a_t = C_0 + ts$, with
$a_{-1} \coloneqq C_0$ by convention. Ordering the classes this way is what makes
prefix masking well defined: it is not a cosmetic relabelling.

Let $f_\theta : \mathcal{X} \to \mathbb{R}^{K}$ be the classifier, emitting a
logit vector of full width $K$ at every stage, and let
$h_\theta : \mathcal{X} \to \mathbb{R}^{d_{\text{lat}}}$ be its penultimate
(latent) representation. Classes that have not yet arrived are suppressed by an
additive mask rather than by reshaping the head,
\begin{equation}
  \label{eq:mask}
  m^{(a)} \in \mathbb{R}^{K}, \qquad
  m^{(a)}_j = \begin{cases} 0, & j < a, \\ -10^{9}, & j \ge a, \end{cases}
\end{equation}
so a single parameterisation carries the whole stream and no logit belonging to
an unseen class can contribute to a loss or to an $\arg\max$.%
\footnote{\texttt{core/training.py:logit\_mask}.}
Table~\ref{tab:notation} collects the remaining symbols.

\begin{table}[t]
  \centering
  \small
  \caption{Notation, and where each quantity is realised in the implementation.}
  \label{tab:notation}
  \begin{tabular}{@{}llp{0.30\linewidth}@{}}
    \toprule
    Symbol & Meaning & Implementation \\
    \midrule
    $\mathcal{D}_t$ & training samples of task $t$ & \texttt{Experiment.task\_loader} \\
    $a_t$ & active logit width after task $t$ & \texttt{TaskSchedule.active\_count} \\
    $m^{(a)}$ & additive prefix mask, Eq.~\eqref{eq:mask} & \texttt{training.logit\_mask} \\
    $\mathcal{M}_t$ & replay buffer after task $t$ & \texttt{core.buffer.ReplayBuffer} \\
    $M$ & total buffer capacity & \texttt{mem\_size} \\
    $\theta^{-}$ & frozen teacher & \texttt{training.clone\_teacher} \\
    $\theta^{\text{pre}}_t$ & model before the forget phase & \texttt{pre} in \texttt{train\_task} \\
    $\mathcal{F}_t,\ \mathcal{R}_t$ & forget / retain split of $\mathcal{D}_t$ & \texttt{select\_forget\_set} \\
    $\rho$ & forget budget as a fraction of $|\mathcal{D}_t|$ & \texttt{forget\_ratio} \\
    $\alpha$ & weight on the forget term & \texttt{alpha} \\
    $\gamma_t$ & weight on the current-task term & \texttt{training.rnt\_for} \\
    $\Omega,\ \theta^{\text{old}}$ & SI importances and anchor & \texttt{SynapticIntelligence} \\
    $c$ & SI strength & \texttt{si\_c} \\
    $\tau$ & distillation temperature & \texttt{kd\_temp} \\
    \bottomrule
  \end{tabular}
\end{table}

\paragraph{Preprocessing.}
Families enter the study only if they carry at least $200$ training samples, and
the $K$ most frequent eligible families are retained. Feature standardisation is
fit \emph{on task~0 alone} and applied to the whole stream, clipped to
$\pm 10$ standard deviations; fitting on the full training set would leak the
statistics of families the learner has not yet met into a protocol whose entire
claim is that it has not met them. LAMDA's features are binary and are left
unscaled.

\subsection{Shared continual backbone}
\label{sec:method:madar}

All conditions share one training loop, one evaluation protocol and one
determinism block, and differ only in the rule applied to a task
(\texttt{experiments/base.py}). We describe that shared backbone first, because
MADAR${+}$Unlearning is defined as a strict extension of it: everything in this
subsection is common to both, so any difference measured between them is
attributable to Section~\ref{sec:method:unlearn} alone.

\paragraph{Task 0.}
The base families are fit by ordinary supervised training for $E_0$ epochs with
Adam,
\begin{equation}
  \label{eq:task0}
  \min_\theta\ \mathbb{E}_{(x,y)\sim\mathcal{D}_0}
    \big[\ \ell_{\mathrm{CE}}\!\big(f_\theta(x) + m^{(a_0)},\, y\big)\ \big].
\end{equation}
Afterwards the SI anchor is set, $\theta^{\text{old}} \leftarrow \theta$, the
teacher is cloned, $\theta^{-} \leftarrow \theta$, and the buffer is filled from
$\mathcal{D}_0$ by the rule of Section~\ref{sec:method:buffer}.

\paragraph{Continual phase.}
For each $t \ge 1$ the model is trained for a \emph{fixed iteration budget}
$N_{\mathrm{cl}}$ rather than a fixed number of epochs, so that compute per task
does not track task size. Each step draws one batch $(x,y) \sim \mathcal{D}_t$
and one replay batch $(\tilde{x},\tilde{y}) \sim \mathcal{M}_{t-1}$ and minimises
\begin{align}
  \label{eq:cl}
  \mathcal{L}^{\mathrm{CL}}_t
   &= \gamma_t\,
      \ell_{\mathrm{CE}}\!\big(f_\theta(x \cup \tilde{x}) + m^{(a_t)},\;
                                y \cup \tilde{y}\big) \nonumber\\
   &\quad + (1-\gamma_t)\,
      \mathcal{L}_{\mathrm{KD}}\!\big(\tilde{x};\, \theta,\, \theta^{-},\, a_{t-1}\big)
      \;+\; c\,\Omega(\theta),
\end{align}
where the distillation term is the temperature-scaled divergence \emph{from} the
teacher \emph{to} the student, restricted to the previous active prefix,
\begin{equation}
  \label{eq:kd}
  \mathcal{L}_{\mathrm{KD}}(x;\theta,\theta^{-},a)
   = \tau^{2}\,
     \mathrm{KL}\!\Big(
       \sigma\big(f_{\theta^{-}}(x)_{:a}/\tau\big)
       \,\Big\|\,
       \sigma\big(f_{\theta}(x)_{:a}/\tau\big)
     \Big),
\end{equation}
with $\sigma$ the softmax. Distillation is taken over \emph{replayed} samples
only, which is why the implementation refuses the configuration
$\texttt{use\_kd}\wedge\neg\texttt{use\_replay}$ outright rather than silently
distilling on the current task.

The convex weight is $\gamma_t = 1/(t+1)$, which weights every task seen so far
equally. Three alternatives (\texttt{floor}, \texttt{sqrt}, \texttt{fixed}) are
implemented and default to off; they exist because $1/(t+1)$ is principled only
when tasks are interchangeable, and under a task~0 far larger than the
increments the newest families receive a small share of the gradient. We note
also that $\gamma_t$ is meaningful only \emph{as} a convex weight against the KD
term: in an ablation that removes distillation, retaining $\gamma_t$ degenerates
into a bare $1/(t+1)$ scalar on the whole objective --- a decaying learning rate
in the costume of a replay weight. It is therefore disabled by default whenever
KD is disabled.

\paragraph{Synaptic Intelligence.}
Parameter importance is the path integral of Zenke et al. During training, each
optimiser step contributes
\begin{equation}
  \label{eq:si-accum}
  W_k \;\mathrel{+}=\; -\,g_k \cdot \Delta\theta_k,
\end{equation}
with $g_k$ the pre-step gradient and $\Delta\theta_k$ the realised parameter
change. At a task boundary the running sum is folded into the importances,
normalised by the squared distance the parameter actually travelled,
\begin{equation}
  \label{eq:si-omega}
  \Omega_k \;\mathrel{+}=\;
    \frac{W_k}{\big(\theta_k - \theta^{\text{old}}_k\big)^{2} + \xi},
  \qquad W_k \leftarrow 0,
\end{equation}
and the quadratic penalty charged during subsequent training is
\begin{equation}
  \label{eq:si-penalty}
  \Omega(\theta) \;=\; \sum_k \Omega_k \big(\theta_k - \theta^{\text{old}}_k\big)^{2}.
\end{equation}
Whether the anchor $\theta^{\text{old}}$ advances at the boundary is a separate
decision from whether $\Omega$ is updated, and Section~\ref{sec:method:order}
turns on exactly that separation.

\paragraph{Frozen normalisation statistics.}
BatchNorm layers are held in evaluation mode for the entire continual phase.
Their running statistics were estimated on task~0, and allowing a stream
dominated by the newest families to overwrite them is a channel of forgetting
that is unrelated to the weights and would otherwise be silently attributed to
the training rule.

\subsection{Diversity-aware replay buffer}
\label{sec:method:buffer}

The buffer holds at most $M$ samples and is budgeted uniformly across the
families seen so far: with $F$ distinct families held or arriving, the
per-family budget is $b = \max(1, \lfloor M / F \rfloor)$, existing families are
truncated to $b$, and each arriving family is filled to $b$.

Within a family, selection is MADAR's diversity criterion. An Isolation Forest
is fitted over the family's samples and they are ordered ascending by
$\mathrm{decision\_function}$, i.e. most anomalous first:
\begin{equation}
  \label{eq:buffer-order}
  \pi = \operatorname{argsort}\big(\mathrm{IF}(\,\cdot\,)\big),
  \qquad
  \mathcal{M}^{(f)} = \operatorname{interleave}\big(
     \pi_{1:\lfloor b/2 \rfloor},\;
     \pi_{n - \lceil b/2 \rceil + 1 : n}\big),
\end{equation}
keeping the most anomalous half and the most prototypical half. The two ends are
\emph{interleaved} rather than concatenated so that a later truncation to a
smaller budget --- which takes a prefix of the stored list --- remains balanced
between the two ends instead of retaining whichever end happens to sit first.

Two spaces are selectable for the forest. Fitting on the scaled \emph{input}
features (\texttt{--buffer-space raw}) is MADAR as published: selection is then a
fixed property of the data. Fitting on the classifier's latent representation
$h_\theta$ (\texttt{latent}, the default here; MADAR-$\theta$) makes selection
move as the model learns. The distinction is not cosmetic --- it decides whether
the memory is chosen in a fixed space or in one the learner is simultaneously
reshaping --- and it is recorded per run.\footnote{The Isolation Forest's
\texttt{contamination} parameter appears in the configuration but only shifts the
estimator's decision \emph{threshold}. Selection here is purely rank-based on
$\mathrm{decision\_function}$ scores, so its value does not change which samples
are chosen; it is retained because the estimator requires it.}

\subsection{MADAR + Unlearning}
\label{sec:method:unlearn}

The proposed condition is the backbone of Section~\ref{sec:method:madar} with a
\emph{forget phase} appended to every task $t \ge 1$. A fraction of the task the
model has just learned is nominated for deletion, the model is trained to become
uninformative about it, and the replay buffer is rebuilt so that the deleted
samples cannot return. The phase runs \emph{after} the continual phase, so the
setting is learn-then-forget: the samples in question have already been fit.
Algorithm~\ref{alg:madaru} states the whole procedure; the subsections that
follow justify the parts that are not mechanical.

\begin{algorithm}[t]
\caption{MADAR${+}$Unlearning, one task $t \ge 1$}
\label{alg:madaru}
\begin{algorithmic}[1]
  \Require model $\theta$, teacher $\theta^{-}$, buffer $\mathcal{M}_{t-1}$,
           SI state $(W,\Omega,\theta^{\text{old}})$
  \State $\theta \gets \textsc{ContinualPhase}(\theta,\theta^{-},\mathcal{D}_t,\mathcal{M}_{t-1})$
         \Comment{Eq.~\eqref{eq:cl}, $N_{\mathrm{cl}}$ steps}
  \State $\Omega \gets \textsc{FoldImportances}(W,\theta,\theta^{\text{old}})$
         \Comment{Eq.~\eqref{eq:si-omega}; anchor \emph{not} advanced}
  \State $(\mathcal{F}_t, \mathcal{R}_t) \gets \textsc{Select}(\mathcal{D}_t, \rho)$
         \Comment{Section~\ref{sec:method:select}}
  \If{$\mathcal{F}_t = \emptyset$}
      \State advance anchor, refresh teacher, rebuild buffer from $\mathcal{D}_t$;
             \textbf{return}
  \EndIf
  \State $\theta^{\text{pre}} \gets \theta$ \Comment{frozen; teacher and reference}
  \For{$e = 1$ \textbf{to} $E_u$}
    \For{minibatch $x_f \subset \mathcal{F}_t$}
      \State draw $x_r \sim \mathcal{R}_t$, \; $\tilde{x} \sim \mathcal{M}_{t-1}$
      \State $\theta \gets \theta - \eta_u \nabla_\theta \mathcal{L}^{U}$
             \Comment{Eq.~\eqref{eq:unlearn}}
    \EndFor
  \EndFor
  \State $\textsc{Measure}(\theta^{\text{pre}}, \theta, \mathcal{F}_t, \mathcal{R}_t)$
         \Comment{Section~\ref{sec:method:efficacy}}
  \State $\theta^{\text{old}} \gets \theta$;\quad $\theta^{-} \gets \theta$
  \State $\mathcal{M}_t \gets \textsc{RebuildBuffer}(\mathcal{M}_{t-1}, \mathcal{R}_t)$
         \Comment{retain split only}
\end{algorithmic}
\end{algorithm}

\subsubsection{Ordering of the SI boundary}
\label{sec:method:order}

The task boundary is closed \emph{before} the forget phase, but the anchor is
deliberately held back: Eq.~\eqref{eq:si-omega} is applied with
$\theta^{\text{old}}$ left pointing at the parameters as they stood at the
\emph{start} of task $t$, and it advances only once unlearning is complete
(lines 2 and 13 of Algorithm~\ref{alg:madaru}).

This ordering is load-bearing in both directions. Folding $W$ into $\Omega$
first means the importances already account for what task $t$ made important, so
the forget phase is constrained by a current estimate rather than a stale one.
Withholding the anchor means the penalty in Eq.~\eqref{eq:si-penalty} is
evaluated against the start of the task, so unlearning is charged for every
important parameter it drags away from that point. Had the anchor advanced
first, $\theta - \theta^{\text{old}} = 0$ at the moment the forget phase begins
and the regulariser would be identically inert for precisely the phase it is
meant to regulate.

We note for completeness that the path integral is \emph{not} accumulated during
the forget phase --- Eq.~\eqref{eq:si-accum} is applied only inside the
continual phase, since $W$ has already been consumed at line~2. The unlearning
trajectory is therefore regularised by SI but does not itself contribute to
future importance estimates.

\subsubsection{Forget-set selection}
\label{sec:method:select}

Selection partitions the current task's samples into
$\mathcal{D}_t = \mathcal{F}_t \sqcup \mathcal{R}_t$ under a budget
$k = \lfloor \rho\,|\mathcal{D}_t| \rfloor$ that is held \emph{fixed across
selectors}. Holding the budget fixed is what makes a comparison between rules
attributable to \emph{which} samples are named rather than to how many. The
seed for the decision is derived, $\mathrm{seed}_t = \mathrm{child}(\mathrm{seed}, t, 7)$,
rather than drawn from the global stream, so that introducing a selector cannot
shift every later random draw in the run and thereby de-pair two conditions
intended to share a seed. If $k = 0$ the phase is skipped and the run record
says so explicitly.

Three rules are compared.

\begin{description}
  \item[\texttt{random}] A uniform subset of size $k$. This is the control for
    \emph{which} samples: it isolates the effect of the selection rule from the
    effect of deleting $k$ samples at all.

  \item[\texttt{donut}] (default) An Isolation Forest is fitted jointly over the
    task in latent space, samples are ordered ascending by
    $\mathrm{decision\_function}$, and the forget set is the \emph{middle band}
    of that ordering,
    \begin{equation}
      \label{eq:donut}
      \mathcal{F}_t = \pi_{\,j:\,j+k},
      \qquad j = \Big\lfloor \tfrac{n}{2} \Big\rfloor - \Big\lfloor \tfrac{k}{2} \Big\rfloor .
    \end{equation}
    Both the most anomalous and the most prototypical samples are kept. This
    applies MADAR's own prior --- the criterion it uses to \emph{choose} what to
    remember, Eq.~\eqref{eq:buffer-order} --- to the complementary question of
    what to delete, and is the hand-crafted rule the learned policy is later
    measured against.

  \item[\texttt{leftover}] A per-family Isolation Forest in \emph{raw} feature
    space; each family keeps the anomalous/prototypical halves that a
    MADAR-style selection at the current per-family budget would have kept, and
    everything else is forgotten. This rule ignores $\rho$ by construction: the
    forget set is whatever the buffer budget leaves over, which is typically far
    larger than a $10\%$ slab. Comparisons involving it are therefore not
    budget-matched, and we report it as such.
\end{description}

\subsubsection{The forget objective}
\label{sec:method:objective}

Let $\theta^{\text{pre}}$ denote the model at the end of the continual phase,
frozen and used both as the distillation teacher for this phase and as the
reference for efficacy measurement. The retain anchor is therefore ``what the
model believed a moment ago'', not the previous task's model. Writing
$u = \mathbf{1}/s$ for the uniform distribution over the $s = a_t - a_{t-1}$
classes this task introduced, the phase minimises
\begin{align}
  \label{eq:unlearn}
  \mathcal{L}^{U}
   &= \alpha\,
      \underbrace{\mathrm{KL}\!\Big(u \,\Big\|\,
        \sigma\big(f_\theta(x_f)_{a_{t-1}:a_t}\big)\Big)}_{\text{forget}}
      \nonumber\\
   &\quad + (1-\alpha)\,
      \underbrace{\Big[
        \tfrac{1}{2}\,\ell_{\mathrm{CE}}\big(f_\theta(x_\cup) + m^{(a_t)},\, y_\cup\big)
        + \tfrac{1}{2}\,\mathcal{L}_{\mathrm{KD}}\big(x_\cup; \theta, \theta^{\text{pre}}, a_t\big)
      \Big]}_{\text{retain}}
      \nonumber\\
   &\quad + c\,\Omega(\theta),
\end{align}
where $x_f \subset \mathcal{F}_t$ and
$(x_\cup, y_\cup) = (x_r,y_r) \cup (\tilde{x},\tilde{y})$ concatenates a batch of
this task's \emph{retained} samples with a replay batch from the buffer. The
forget set defines the step count; the retain and replay streams cycle to match
it. Optimisation uses Adam at $\eta_u$ for $E_u$ passes over $\mathcal{F}_t$,
with gradient-norm clipping at $0.5$ and BatchNorm still frozen.

Three properties of Eq.~\eqref{eq:unlearn} should be read carefully.

\emph{The forget term acts on the new-task slice only.} The softmax is taken over
$[a_{t-1}, a_t)$ --- the $s$ classes this task introduced --- not over the full
active prefix. The objective therefore drives the model toward indifference
\emph{among the new families}, which is a narrower claim than indifference among
all classes seen so far.

\emph{Retention is anchored on two populations.} Because $x_\cup$ mixes retained
current-task samples with replayed old-family samples, the $(1-\alpha)$ term
defends both plasticity on the task being partially forgotten and stability on
everything before it, under one weight.

\emph{The regulariser is not slack.} By Section~\ref{sec:method:order}, $\Omega$
is evaluated against the start of the task, so the forget term must overcome an
active quadratic penalty on exactly those parameters the task made important.

\subsubsection{Buffer reconstruction from the retain split}
\label{sec:method:rebuild}

After the forget phase the buffer is rebuilt from $\mathcal{R}_t$ rather than
$\mathcal{D}_t$: latents are recomputed with the \emph{post}-unlearning model and
Eq.~\eqref{eq:buffer-order} is re-applied per family. Forgotten samples are
consequently never eligible for replay at any later task.

This is worth separating from the objective, because it is a second and
independent deletion mechanism, and it operates whenever a forget set is selected
--- including when $\alpha = 0$. It is the reason the ablation of
Section~\ref{sec:method:controls} is necessary rather than decorative.

\subsubsection{Efficacy measurement}
\label{sec:method:efficacy}

Unlearning is measured, per task, by comparing $\theta^{\text{pre}}$ against the
post-phase $\theta$ on the same inputs, over two disjoint scopes: the forget set
$\mathcal{F}_t$ and the retained current-task samples $\mathcal{R}_t$. For each
scope we record accuracy before and after, the mean predictive entropy and its
normalised form, and the displacement of the latent representation:
\begin{align}
  \label{eq:efficacy}
  H(x) &= -\!\!\sum_{j < a_t} p_j \log p_j,
  \qquad p = \sigma\big(f_{\theta}(x)_{:a_t}\big), \\
  \widehat{H} &= \frac{\mathbb{E}[H(x)]}{\log a_t},
  \qquad
  \Delta_{\text{lat}}(x) = \big\| h_{\theta}(x) - h_{\theta^{\text{pre}}}(x) \big\|_2 .
\end{align}
A useful result is a large drop in accuracy and a large $\Delta_{\text{lat}}$ on
$\mathcal{F}_t$ together with small movement on $\mathcal{R}_t$: forgetting that
also flattens the retained samples is collateral damage, not deletion. Note that
$\widehat{H}$ is normalised over the \emph{full} active prefix while the forget
term of Eq.~\eqref{eq:unlearn} operates on the $s$-class slice, so a forget set
driven to perfect uniformity within that slice will not register
$\widehat{H} \approx 1$; the two widths are not interchangeable and we report
$\widehat{H}$ as a monotone diagnostic rather than as a saturating one.

\subsection{Conditions, controls and ablations}
\label{sec:method:controls}

Four conditions share the network, the family selection, the scaler, the task
schedule, the masked evaluation and the determinism block. Only the rule applied
to a task differs.

\begin{table}[t]
  \centering
  \small
  \caption{The four conditions. All share every non-training component.}
  \label{tab:conditions}
  \begin{tabular}{@{}lp{0.62\linewidth}@{}}
    \toprule
    Condition & Training rule \\
    \midrule
    Naive & Fine-tune on the current task only. Lower bound. \\
    Joint & Discard the model and retrain from scratch on all data seen so far.
            Oracle upper bound; not a realisable continual setup. \\
    MADAR & Eq.~\eqref{eq:cl}: diversity-aware replay, distillation against the
            previous task's model, SI penalty. \\
    MADAR${+}$U & MADAR, plus the per-task forget phase of
            Section~\ref{sec:method:unlearn}. \\
    \bottomrule
  \end{tabular}
\end{table}

\paragraph{$\alpha = 0$ is the matched control.}
The headline comparison is not MADAR${+}$U against MADAR. At $\alpha = 0$ the
forget set is still selected by the same rule, is still excluded from the buffer
rebuild, and the same number of gradient steps is still taken --- only the push
toward uniform is absent. The contrast $\alpha > 0$ versus $\alpha = 0$ therefore
isolates the value of the \emph{forget objective}, while $\alpha = 0$ versus
MADAR isolates the value of \emph{curation}, i.e. of removing those samples from
memory at all. Reporting the method without the $\alpha = 0$ arm would attribute
to unlearning what may simply be memory curation, and both arms are run for every
corpus and seed.

\paragraph{Component ablations.}
Replay, distillation and the SI penalty are independently switchable in
Eq.~\eqref{eq:cl} so that an ablation removes exactly one thing. SI is ablated by
setting $c = 0$, which leaves the optimiser, the iteration count and the data
order untouched; removing distillation additionally disables $\gamma_t$ for the
reason given in Section~\ref{sec:method:madar}.

\subsection{Evaluation protocol}
\label{sec:method:eval}

After every task, predictions are taken as
$\hat{y} = \arg\max_{j < a_t} f_\theta(x)_j$ over the test split, and all
reported quantities are derived from a single confusion matrix so that the
numbers in a record are mutually consistent by construction rather than by three
code paths agreeing. Three scopes are evaluated:
\begin{itemize}
  \item \textbf{seen} --- all families encountered so far. The headline number.
  \item \textbf{recent} --- the families just introduced. Acquisition, i.e.
        plasticity.
  \item \textbf{task0} --- the base families. Retention, i.e. resistance to
        forgetting.
\end{itemize}
The \textbf{seen} scope averages the other two, so a method that trades
acquisition against retention is invisible in it alone; we report all three.

Macro averages are taken over classes \emph{with support in the evaluation
slice}. A class absent from a slice has no recall to measure, and scoring it as
zero would make the metric a function of which classes happen to be absent.
Because a class may be predicted without being present, the false positives such
predictions generate belong to no scored class; they are surfaced explicitly as
\texttt{n\_pred\_outside\_support} rather than dropped. For single-label
multiclass, micro-precision, micro-recall and micro-F1 all equal accuracy, so
only accuracy is reported under that name.

Two protocol constraints are enforced rather than documented. First, the
\texttt{stage1} protocol (task~0 retrained per seed) and the \texttt{meta}
protocol (a shared task-0 checkpoint) are never differenced against each other,
and are tabulated as separate blocks. Second, a crashed run is recorded as failed
and is never scored, defaulted or interpolated --- a missing number that silently
becomes \texttt{NaN} and then sorts to the top of a ranking is a real failure
mode, and the sweep boundary is where it is refused.

\subsection{Hyperparameters}
\label{sec:method:hparams}

Table~\ref{tab:hparams} gives the values in effect. Per-corpus entries are
protocol for that corpus; an explicit override always wins and, whatever wins,
the resolved value is written into the run record --- there is no value in effect
that the log does not name.

\begin{table}[t]
  \centering
  \small
  \caption{Protocol parameters. Shared values are common to every corpus.}
  \label{tab:hparams}
  \begin{tabular}{@{}llcccc@{}}
    \toprule
    & & EMBER 2018 & EMBER 2024 & LAMDA & Tiny ImageNet \\
    \midrule
    \multicolumn{6}{@{}l}{\emph{Shared}}\\
    Task-0 epochs $E_0$        & \texttt{task0\_epochs}   & \multicolumn{3}{c}{30, Adam $10^{-3}$} & 50 \\
    Batch size                 & \texttt{batch\_size}     & \multicolumn{3}{c}{256} & 32 \\
    Continual optimiser        & ---                      & \multicolumn{4}{c}{SGD, momentum $0.9$, weight decay $10^{-6}$, clip $1.0$} \\
    Distillation temperature   & \texttt{kd\_temp}        & \multicolumn{4}{c}{$\tau = 2.0$} \\
    Current-task weight        & \texttt{rnt\_mode}       & \multicolumn{4}{c}{$\gamma_t = 1/(t{+}1)$} \\
    Forget weight              & \texttt{alpha}           & \multicolumn{4}{c}{$\alpha \in \{0.2,\ 0\}$ (both arms)} \\
    Forget budget              & \texttt{forget\_ratio}   & \multicolumn{4}{c}{$\rho = 0.10$} \\
    Forget passes              & \texttt{unlearn\_epochs} & \multicolumn{3}{c}{$E_u = 3$} & 3 \\
    Forget optimiser           & \texttt{unlearn\_lr}     & \multicolumn{4}{c}{Adam $\eta_u = 10^{-4}$, clip $0.5$} \\
    Family eligibility         & \texttt{min\_family\_samples} & \multicolumn{3}{c}{$\ge 200$} & --- \\
    \midrule
    \multicolumn{6}{@{}l}{\emph{Per corpus}}\\
    Schedule                   & \texttt{--task-setup}    & $50{+}5{\times}10$ & $50{+}5{\times}10$ & $30{+}5{\times}10$ & $100{+}10{\times}10$ \\
    Continual budget $N_{\mathrm{cl}}$ & \texttt{cl\_iters} & 2000 & 8000 & 4000 & 3000 \\
    Continual lr               & \texttt{cl\_lr}          & $10^{-4}$ & $10^{-4}$ & $10^{-4}$ & $0.03$ \\
    Buffer capacity $M$        & \texttt{mem\_size}       & 5000 & 10000 & 2500 & 2000 \\
    SI strength $c$            & \texttt{si\_c}           & 1.0 & 100.0 & 100.0 & 1.0 \\
    Per-family cap             & \texttt{family\_cap}     & --- & 10000 & --- & --- \\
    Feature clip               & \texttt{feature\_clip}   & $\pm10$ & $\pm10$ & --- (binary) & --- \\
    \bottomrule
  \end{tabular}
\end{table}

\subsection{Threats to validity}
\label{sec:method:threats}

\paragraph{Compute is not matched between MADAR and MADAR${+}$U.}
The forget phase adds $E_u \lceil |\mathcal{F}_t| / B \rceil$ optimiser steps per
task on top of the $N_{\mathrm{cl}}$ continual steps, and the run record reports
the total. The difference is small relative to $N_{\mathrm{cl}}$ but it is not
zero, so a MADAR${+}$U improvement over MADAR is confounded with additional
gradient steps. The $\alpha = 0$ arm of Section~\ref{sec:method:controls} is
matched in step count and is the comparison from which the forget objective's
contribution should be read.

\paragraph{The forget objective and the efficacy metric span different widths.}
As noted in Section~\ref{sec:method:efficacy}, $\mathcal{L}^U$ flattens the
$s$-class new-task slice while $\widehat{H}$ normalises over the full active
prefix of width $a_t$. The two are consistent in direction but not on a common
scale.

\paragraph{Unlearning is behavioural, not certified.}
Eq.~\eqref{eq:unlearn} drives the model's \emph{predictions} on $\mathcal{F}_t$
toward uniformity and removes those samples from memory; it offers no guarantee
of indistinguishability from a model retrained without them, and we make no such
claim. The reported quantities of Section~\ref{sec:method:efficacy} are evidence
of behavioural forgetting under a specific probe, and the \texttt{Joint} oracle
of Table~\ref{tab:conditions} is a retraining reference for accuracy, not a
privacy reference.

\paragraph{The \texttt{leftover} selector is not budget-matched.}
Its forget set is determined by the buffer budget rather than by $\rho$, so
differences against \texttt{random} and \texttt{donut} confound the selection
rule with the number of samples deleted.

\paragraph{Reproduction.}
Every reported run is produced by a single command and writes one self-describing
JSON record naming the corpus, the schedule, the model and parameter count, every
hyperparameter in effect, the seed, the protocol, the library versions, the GPU
and the git commit:
\begin{verbatim}
./iclr run --dataset ember2018 --task-setup ember \
           --experiment madar_unlearn --selector donut \
           --model full --seed 1 --group <group>
\end{verbatim}
